%  LaTeX support: latex@mdpi.com 
%  For support, please attach all files needed for compiling as well as the log file, and specify your operating system, LaTeX version, and LaTeX editor.

%=================================================================
\documentclass[jimaging,article,accept,pdftex,moreauthors]{Definitions/mdpi} 

\graphicspath{{figures/}} % path to folder with figures
\usepackage{caption}
 \usepackage[labelformat=simple]{subcaption}
 \renewcommand\thesubfigure{\alph{subfigure}}
\DeclareCaptionLabelFormat{subcaptionlabel}{\normalfont(\textbf{#2}\normalfont)}
\captionsetup[subfigure]{labelformat=subcaptionlabel}
%\usepackage[export]{adjustbox}
\usepackage{listings}
\usepackage{xcolor}
\usepackage{graphicx}
\usepackage{gensymb} % for degree symbol
\usepackage{lscape}
\usepackage{pxfonts}

\definecolor{codegreen}{rgb}{0,0.6,0}
\definecolor{codegray}{rgb}{0.5,0.5,0.5}
\definecolor{codepurple}{rgb}{0.58,0,0.82}
\definecolor{backcolour}{RGB}{255,255,255}
\definecolor{SlateGray}{HTML}{708090}

\lstdefinestyle{mystyle}{
    backgroundcolor=\color{backcolour},   
    keywordstyle=\bfseries,
    numberstyle=\tiny\color{codegray},
    stringstyle=\color{SlateGray},
    basicstyle=\ttfamily\footnotesize,
%    basicstyle=\ttfamily\scriptsize,
    breakatwhitespace=false,         
    breaklines=true,                 
    captionpos=t, 
    keepspaces=true,                 
    numbers=left,                    
    numbersep=5pt,                  
    showspaces=false,                
    showstringspaces=false,
    showtabs=false,                  
    tabsize=2,
    frame=topline|bottomline,
    }

\lstset{style=mystyle}
\captionsetup[lstlisting]{position=top, margin=0cm,labelfont={bf, small, stretch=1.17}, labelsep=colon, textfont={small, stretch=1.17}, aboveskip=6pt, singlelinecheck=off, justification=justified}
	
\usepackage{soul} % highlighting text

%--------------------
% Class Options:
%--------------------
%----------
% journal
%----------

%---------
% article
%---------

%=================================================================
% MDPI internal commands
\firstpage{1} 
\makeatletter 
\setcounter{page}{\@firstpage} 
\makeatother
\pubvolume{1}
\issuenum{1}
\articlenumber{0}
\pubyear{2022}
\copyrightyear{2022}
\externaleditor{Academic Editor: 
}
\datereceived{14 October 2022} 
%\daterevised{} % Only for the journal Acoustics
\dateaccepted{20 November 2022} 
\datepublished{} 
%\datecorrected{} % Corrected papers include a "Corrected: XXX" date in the original paper.
%\dateretracted{} % Corrected papers include a "Retracted: XXX" date in the original paper.
\hreflink{https://doi.org/} % If needed use \linebreak
%\doinum{}
%------------------------------------------------------------------
% The following line should be uncommented if the LaTeX file is uploaded to arXiv.org
%\pdfoutput=1

%=================================================================
% Add packages and commands here. The following packages are loaded in our class file: fontenc, inputenc, calc, indentfirst, fancyhdr, graphicx, epstopdf, lastpage, ifthen, lineno, float, amsmath, setspace, enumitem, mathpazo, booktabs, titlesec, etoolbox, tabto, xcolor, soul, multirow, microtype, tikz, totcount, changepage, attrib, upgreek, cleveref, amsthm, hyphenat, natbib, hyperref, footmisc, url, geometry, newfloat, caption

%=================================================================
% Full title of the paper (Capitalized)
\Title{Satellite Image Processing by Python and R Using Landsat 9 OLI/TIRS and SRTM DEM Data on Côte d'Ivoire, West Africa}

% MDPI internal command: Title for citation in the left column
\TitleCitation{Satellite Image Processing by Python and R Using Landsat 9 OLI/TIRS and SRTM DEM Data on Côte d'Ivoire, West Africa}

% Author Orchid ID: enter ID or remove command
\newcommand{\orcidauthorA}{0000-0002-5759-1089} % Add \orcidA{} behind the author's name
\newcommand{\orcidauthorB}{0000-0002-6461-1551} % Add \orcidB{} behind the author's name

% Authors, for the paper (add full first names)
\Author{Polina Lemenkova %MDPI: Please carefully check the accuracy of names and affiliations. <--  Answer: yes, correct.
*\orcidA{} and Olivier Debeir \orcidB{}}
% \dagger,\ddagger
%\longauthorlist{yes}

% MDPI internal command: Authors, for metadata in PDF
\AuthorNames{Polina Lemenkova $^{1}$ and Olivier Debeir $^{1}$}

% MDPI internal command: Authors, for citation in the left column
\AuthorCitation{Lemenkova, P.; %MDPI: Please carefully check the accuracy of names. <--  Answer: yes, correct.
 Debeir, O.}
% If this is a Chicago style journal: Lastname, Firstname, Firstname Lastname, and Firstname Lastname.

% Affiliations / Addresses (Add [1] after \address if there is only one affiliation.)
\address[1]{Laboratory of Image Synthesis and Analysis (LISA), École Polytechnique de Bruxelles (EPB, Brussels Faculty of Engineering), Université Libre de Bruxelles. %MDPI: please confirm address order according to order suggestion (no order requirement within the same pair of parentheses): (Laboratory | Group | Programme | Unit) < (Department | Division | Faculty | Institute | School) < (College | University | University) <-- Answer: corrected according to order suggestion starting from Laboratory to University.
Building L, Campus de Solbosch, ULB--LISA CP165/57, Avenue~Franklin D. Roosevelt 50, B-1050 Brussels, Belgium; olivier.debeir@ulb.be %MDPI: We added the email addresses here according to those submitted online at susy.mdpi.com. Please confirm. <--  Answer: yes, correct.
}

% Contact information of the corresponding author
 \corres{\hangafter=1 \hangindent=1.05em \hspace{-0.82em} Correspondence: polina.lemenkova@ulb.be; Tel.: +32-471860459
 }

% Abstract
\abstract{In this paper, we propose an advanced scripting approach {using Python and R for satellite image processing and modelling terrain in} Côte d'Ivoire, West Africa. {Data include Landsat 9 OLI/TIRS C2 L1 and the SRTM digital elevation model (DEM)}. {The EarthPy library of Python and `raster' and `terra' packages of R are used as tools for data processing}. The methodology includes computing vegetation indices {to derive information on vegetation coverage} and {terrain} modelling. Four vegetation indices {were} computed and visualised using R: the Normalized Difference Vegetation Index (NDVI), Enhanced Vegetation Index 2 (EVI2), Soil-Adjusted Vegetation Index (SAVI) and Atmospherically Resistant Vegetation Index 2 (ARVI2). The SAVI index is demonstrated to be more suitable and better adjusted to the vegetation analysis, which is beneficial for agricultural monitoring in C\^ote d'Ivoire. The terrain analysis {is} performed using Python and includes slope, aspect, hillshade and relief modelling {with changed parameters for the sun azimuth and angle}.~The vegetation pattern in Côte d'Ivoire is heterogeneous, which reflects the complexity of the terrain structure. Therefore, {the} terrain and vegetation data modelling {is aimed at the analysis of} the relationship between the regional topography and environmental setting in the study area.~{The upscaled mapping is performed as} regional environmental analysis of the Yamoussoukro surroundings and local topographic modelling of the Kossou Lake.~{The} algorithms of the data processing include image resampling, band composition, statistical analysis and map algebra used for calculation of the vegetation indices in Côte d'Ivoire. This study {demonstrates the effective application of} the advanced programming algorithms {in Python and R for satellite } image processing.}

% Keywords
\keyword{image processing; remote sensing; programming; Python; R; satellite image; digital terrain model; cartography; mapping; spatial analysis} 

\PACS{91.10.Da; 91.10.Jf; 91.10.Sp; 91.10.Xa; 96.25.Vt; 91.10.Fc; 95.40.+s; 95.75.Qr; 95.75.Rs; 42.68.Wt %MDPI: please confirm if these three parts need keep. <--  Answer: yes, correct. PACS, MSC and JEL are needed for indexation of the paper.
}
\MSC{86A30; 86-08; 86A99; 86A04}
\JEL{Y91; Q20; Q24; Q23; Q3; Q01; R11; O44; O13; Q5; Q51; Q55; N57; C6; C61}

%%%%%%%%%%%%%%%%%%%%%%%%%%%%%%%%%%%%%%%%%%
\begin{document}

\section{Introduction}

%---------------------%
\subsection{Background}

Remote sensing (RS) data have great potential {for geosciences as an} important source of spatial information. {The RS data can be used in a variety of applications, including} geologic {studies, land cover mapping,} climate {change detection} and {environmental monitoring}. Semantic information obtained from {satellite images is widely} used in thematic mapping for analysis of {the} diverse phenomena and processes on Earth{. Examples include mapping} land cover fractions \cite{COULIBALY2021108092} and vegetation types \cite{7729337,Lemenkova2020d,6946534,Lemenkova202019,6352694}, hydrologic and geomorphic modelling \cite{2015hydrologie},{ }urban sprawl {and} geological exploration{ }\cite{1457509,NGOM2022102873}. Aside from satellite images, another source of {RS data }is provided by airborne {images} \cite{JESSELL2015440,SIAGNE2022104680} and light detection and ranging (LiDAR) data for three-dimensional (3D) terrain analysis \cite{SHARMA2021317,8978301,9213897,YEH2017236,9553919,LI201955,GUPTA2022100817}. 

The general idea {in image} processing is that{ geospatial information} can be {derived from the interpreted RS data}. {For instance, }repeated structures representing similar objects on the Earth's surface are commonly depicted on satellite images, represented by pixels comprising {the} matrix of the image scene with a comparable range of values. These data can be{ }visualised using the properties of the digital numbers (DN) of the pixels, which depend on the texture, composition and fabric of the objects shaping the face of the Earth \cite{YUAN2022106966}.{ Different} land cover types which comprise the mosaic of the diverse landscapes can be detected {on the images}. Identifying and recognising these features is possible using the {contrasting }parameters of pixels in the satellite images: brightness, resolution \cite{4426419,322052,5446016}, intensity \cite{4761020,8519176} and texture \cite{995559,899813,8868688}.~{Accurate and robust }techniques{ }of satellite image processing {are} essential {for effective image interpretation in many} approaches{ }developed {recently}. These include{ }diverse algorithms of image classification and processing{.} 

{The existing }algorithms of image classification{ aim at reading the information from pixels through reading the }values{ of} raster cells. The analysis of this information enables mapping{ }diverse phenomena, processes and objects based on {the} interpretation of {RS} data \cite{PERIN2021106694,PAHLEVAN2022112860,Lemenkova202141,ABU2021107863,ATTOUMANE2022103285}. The main objective of image processing is {obtaining} this information from the {RS} data{. }In this context, the search for effective methods of image processing resulted in the variety of existing Geographic Information System (GIS) software{ designed }for {RS} data processing. However, in many practical situations, the GIS presents limitations, either due to the commercial licencing of software {with} an associated high cost {(}ArcGIS, Erdas Imagine and eCognition) or functional restrictions (QGIS and the System for Automated Geoscientific Analyses (SAGA) GIS). Moreover, the {traditional manual }methods{ of image processing are time-consuming, which make them ineffective for big} data processing {and require automation}. 

%---------------------%
\subsection{Related Work}

Programming {languages} present{ }effective solutions {for processing} satellite images{. This is achieved through the} advanced methods of machine learning (ML) \cite{MASOLELE2021112600,9066075,MARTINS2022113203,rs71114876}. Among them, scripting techniques enable handling {RS} data effectively and accurately{ }\cite{6344536,ELLEFSEN20191171,Lemenkova202212,6910592,Lemenkova202215,4160259}. {An }essential advantage of the programming languages consists of {automation of the RS data processing. A second advantage consists of the }compatibility{ of data formats}. In contrast to the stand-alone GIS software, scripting libraries can be integrated with other GISs via modules, plugins and toolkits \cite{6407912,7378577,VOS2019104528}. A variety of existing programming languages {can be used as} optimal and effective tools{ }for image processing and {RS} data analysis in environmental modelling,{ among which Python and R are the most advanced and widely used.}

Python is undoubtedly one of the most successful high-level programming languages~\cite{van1995python}. It has become one of the most powerful tools that can be used for data science. {Python} enables performing various tasks of image processing and spatial analysis through integration {of libraries} \cite{1247428,Rey2010PySALAP,Rey2010}. Recent examples include image interpolation, calibration and correlation \cite{DEDEUSFILHO2022101131}, compressed sensing and image reconstruction \cite{FARRENS2020100402}, identifying and extracting vertical features \cite{GAO2022105503} and meteorological modelling \cite{6528254} with resulting images in compatible formats. 

A collection of Python libraries \cite{ROBERTS2022105192,CHEN2022105362} has recently built up with EarthPy that is specially tailored for{ }geo-information processing \cite{wasser2019EarthPyAP}.~Aside from general applications in {diverse branches of }data science and ML \cite{8757088,DebeirDecaestecker,Debeir2008,9719489}, Python can be applied for {RS data processing using its spatial libraries} \cite{Ahmed}. Image and signal processing by {its} object-oriented algorithms of Python are pervasive in geosciences and provide solutions to a wide range of problems{. }Examples of using Python have been reported {in a variety of domains, including} geodetic studies \cite{Zhou2021GPS2spaceAO}, topographic analysis \cite{9254627}, hydrogeological modelling \cite{Mogaji}, modelling air pollution and land use \cite{Maetal2020}, batch spatial data processing \cite{6057428}, creating panorama images~\cite{JASKOLKA2019e02560} and tunneling microscopy \cite{BASTIDAS2021150821}, {to mention a few}. 

{R }is a fundamental programming tool for {statistical and }spatial data processing \cite{RCoreTeam}. It { includes a} variety of packages{ with extended} graphical functionality{ and advanced} statistical components \cite{Murrell}.~The packages of R have long been used for data retrieval, analysis, modelling and visualisation in various industries including robotics, engineering and GIS \cite{7847824,Lemenkova2022a,9378399,Lemenkova202216,9390011}. R demonstrates{ }superiority in statistical computing, {graphical plotting and data analysis compared with} graphical user interface (GUI) software. {Moreover, R} excels in big data analysis \cite{7570960,9752002,6597171}, data mining \cite{7978404} and visualisation \cite{Lemenkova201991}{ }modelling, plotting{ and image processing, which are }supported through the variety of its built-in packages. 

{Compared} with Python, R has a{ }richer collection of libraries{ due to the }long history of stable development{. They contain} a wide variety of mathematical statistical tools, from the very basic functions to the highly sophisticated algorithms of data analysis~\cite{9675988}. {Recent} advances in geospatial data processing \cite{8226275} and image classification~\cite{9844952,8117873} have made R successful for remote sensing data processing, including image classification and segmentation. This poses a new way to process and model satellite images and derive spatially structured information from complex scenes. {In this regard, the 'terra' and 'raster' libraries of R present essential tools for }modelling geospatial {data.~The functionality of these libraries includes essential options for image classification, such as reading spatial dimensions from raster images, supporting spectral bands of Landsat scenes, evaluating the} similarity between {the values of pixels and extracting features from variables such as} vegetation types and variations in landscape structure and topography. 

%---------------------%
\subsection{Contribution}

{The }goal of this paper is to present the application of Python and R algorithms for image processing{. This work is motivated by }the idea of using{ }advanced programming methods{ }to visualise{ }object features and classes on the Earth's surface{. The }specific objectives include the following tasks: (1) %MDPI: Please confirm if the italics is unnecessary and can be removed. The following highlights are the same. <--  Answer: yes, can be removed. 
plotting the colour composite of a Landsat Operational Land Imager and Thermal Infrared (OLI/TIRS) image by{ } combinations of {the} three channel sensor bands to display the image; (2) computing the selected vegetation indices using the principles of map algebra based on {the} DN of reflectance value of the pixels in the selected spectral bands of near infrared (NIR) and red; and (3) terrain modelling with geomorphic visualisation of the slope, aspect, hillshade and elevation for relief analysis using the {data derived from the }Shuttle Radar Topography Mission (SRTM) digital elevation model (DEM). 

{We propose} two approaches{ }to{ }satellite image processing by{ }Python and R{. }In the first approach, we {use }R packages `terra' and `raster' for calculation of the vegetation indices {of Landsat bands} and creating colour band composites.~For the second approach, we {present} an application of the Python library EarthPy for terrain analysis supported by auxiliary libraries (Matplotlib, NumPy, Rasterio and Plotly) for data processing and visualisation. The combination of several libraries {has been demonstrated to be effective in utilising} the properties and functionality of each to achieve effective image processing for {geospatial} analysis {using the Python syntax}.

The objectives of our study {are} to use the advanced functionality of{ }R and Python for {satellite} image {analysis }aimed at {the} environmental assessment and mapping of C\^ote d'Ivoire{ using broadband multispectral images. Here, the goals are to identify the vegetation conditions by computing four vegetation indices---NDVI, SAVI, EVI2 and ARVI2---and to model} the terrain in the surroundings of Kossou Lake to the north of Yamoussoukro. {Several libraries are used in scripts for }terrain modelling, computing the vegetation indices {and creating colour composites using a} combination of Landsat bands. {In contrast to the existing approaches, which use the traditional GIS software with restricted functionality,} our method {applies} flexible algorithms of scripting languages {using the} syntax {of Python and R and extended functionality of their libraries}. 

{We find that} the advanced libraries of Python and{ R are effective} for processing high-resolution satellite images{. This motivates us to introduce an efficient scripting method to environmental monitoring. }To the best of our knowledge, no attempts have been made to {integrate} high-level{ }languages, such as Python and R, for mapping C\^ote d'Ivoire, {for} terrain modelling or for vegetation {analysis by scripts}. Works related to ours {include} some previous {studies on C\^ote d'Ivoire} \cite{Chatelainetal2010,Hennenberg,ElShahat,Tangetal2016,Affianetal2009}, in which traditional GIS methods were applied to process spatiotemporal information for environmental mapping{. }The key difference between our work and theirs is that our method integrates programming algorithms {for processing multispectral} satellite images {using} scripts.

Our main contributions can be summarised as follows: (1){ }environmental mapping of C\^ote d'Ivoire based on modelling and visualisation of several vegetation indices using the `terra' package; (2){ }an R-based approach to processing satellite images{ }for automatic processing of raster data in spectral bands and computing cell values {aimed at} spatial analysis of the heterogeneous landscapes {of C\^ote d'Ivoire}; (3){ }EarthPy-based algorithms {for deriving} topographic information of the Kossou Lake region{ to} model the terrain parameters of the slope, aspect, hillslopes and stream channels of the hydrological network system; and (4) mapping spatial data with scripting libraries {for automatic data processing,} which {outperforms the traditional approaches} for cartographic visualisation {and image analysis}.  

\section{Study Area}

The study area was focused on the central part of Côte d'Ivoire, West Africa (Figure \ref{fig01}). In particular, {we selected} the region of Kossou Lake{ }as an enlarged fragment {of the study area} for{ }terrain analysis. This lake {was formed} artificially {by a constructed} dam {on} the Bandama River. {Kossou} Lake forms one of the most important inland water bodies of Côte d'Ivoire and plays an essential{ }role in the environment of Sub-Saharan Africa \cite{MSISKA2009487}. It has functional aspects{ }as a habitat for dominant biotas and {supports} food {chains in ecosystems}. 
\vspace{-12pt}
\begin{figure}[H]

	\includegraphics[width=11.0 cm]{Fig_01}
	\caption{Topographic map of Côte d'Ivoire. Mapping software: Generic Mapping Tools (GMT) scripting toolset. Data source: GEBCO/SRTM. Cartography: authors.
	\label{fig01}}
\end{figure} 

C\^ote d'Ivoire is notable for {a} diverse terrain structure which has two distinct parts: a coastal region formed by the alluvium and marine sands and a continental region with hilly relief \cite{Thomasset}. The land cover types are associated with two major ecosystem types: forests{ }distributed{ }in the southern part of the country{ }and savannah to the north off Bouak\'{e} \cite{Rougerie}. The vegetation types are largely controlled by the equatorial or subequatorial rainfall regime.~A long dry season is typical for the northern half of Côte d'Ivoire with dominating savannahs, while a tropical regime characterises the southern part of the country with prevailing forests~\cite{Tricart}.~A favourable climatic and environmental setting maintains the farming and agricultural activities {of the local population. This makes} Côte d'Ivoire the largest cocoa producer in the world \cite{Losch,Planhol,Smithetal2014,Laderach} and an active exporter of food-producing crops \cite{Sawadogo}, such as coffee, bananas, pineapples and palm oil \cite{Pelissier}. Aside from agriculture, the activities of population include exploration, mining \cite{SAKO2018297,SAUERWEIN2020104323,MURRAY2019351} and limited fishing along the narrow continental shelf of the Atlantic Ocean \cite{Sournia}. 

Recently, Côte d'Ivoire {experienced} environmental changes{ typical for other countries of} humid West Africa{. These }include deforestation and fragmentation of ecosystems~\cite{Chatelain}, a decrease in species richness in the savannah biomes \cite{Kolongo} and an increase in habitat heterogeneity, which affects biodiversity and the species community structure \cite{Yeo}. {The triggers for the environmental changes include} climate effects {and} urban sprawl, {which} has been exceptional for Côte d'Ivoire since the 1950s, due to the natural demographic growth and the {increase in} industrial agglomerations \cite{Dubresson,Chaleard,Cotten1974,Cotten1973}.~Geographically, {the} acceleration of the population {leads to} increased city areas and modified transport networks \cite{Cotten1985}. Other environmental problems include water supply issues and pollution caused by agro-industrial plantations \cite{Doumouya}. 

%%%%%%%%%%%%%%%%%%%%%%%%%%%%%%%%%%%%%%%%%%
%\newpage
\section{Materials and Methods}

\subsection{Data}
The data were collected from the United States Geological Survey (USGS)'s Global Visualization Viewer (GloVis), a satellite image repository with {publicly available} open access to remotely sensed data (Figure \ref{fig02}). 

\begin{figure}[H]
     \begin{subfigure}[r]{0.60\textwidth}
         
         \includegraphics[width=9.5cm]{Fig_02a}
         \caption{\centering}
     \end{subfigure} \vfill \vspace{2mm}
     \begin{subfigure}[r]{0.60\textwidth}
         \includegraphics[width=9.5cm]{Fig_02b}
         \caption{\centering}
     \end{subfigure}
        \caption{Data captured from the USGS Global Visualization Viewer (GloVis)'s satellite image repository, accessing remote sensing data from Landsat 9 OLI/TIRS C2 L1 and SRTM DEM for terrain analysis. ({\bf a})~Collected from Landsat 9 OLI/TIRS C2 L1. %MDPI: We moved the subfigure explanations into the figure caption. Please confirm. <--  Answer: yes, agree.
 ({\bf b}) Collected from SRTM void-filled DEM.}
        \label{fig02}
\end{figure}

The {data} used {in this study included the satellite imaging of Landsat 9 OLI/TIRS C2 L1 and the SRTM DEM for terrain modelling. The} satellite image {covers the} central part of Côte d'Ivoire{. It has the following technical parameters: }a LC91970552022011LGN01 scene identified; Landsat product ID of LC09\_L1TP\_197055\_20220111\_20220122\_02\_T1 and land cloud coverage of 1.29, acquired on 11 January 2022. The{ }image{ was }obtained from the Landsat 9 Earth observation satellite, which has two remote sensing instruments {for} optical and thermal sensors onboard: the Operational Land Imager (OLI-2) {sensor} and {the} Thermal Infrared Sensor (TIRS-2) along with nine spectral bands \cite{USGS,Landsat9}. More technical information is available online from {Landsat 9}. The data used for terrain analysis {included the} SRTM void-filled DEM, {obtained} from the {SRTM3 NASA/NGA}{. It has a resolution of} 3 arc-seconds{ (90 m) and the following technical parameters:} an entity ID of SRTM3N07W006V2 and tile coordinates: of 5°--6° W, 7°--8° N{, with} an acquisition {date} of 11 February 2000. The general map showing the study area {was} plotted based on the General Bathymetric Chart of the Oceans (GEBCO) \cite{GEBCO} {using} GMT scripting \cite{Wessel} and {following the} existing methodology \cite{Lemenkova2022b,Lemenkova202217}. 

%\newpage
\subsection{Image Processing in R}
The R packages `terra' \cite{RTerra} and `raster' \cite{Hijmans} were used for computing the vegetation indices and image processing. The colour scheme resource of the RColorBrewer package was used for visualisation and plotting \cite{Neuwirth}.

\subsubsection{Information Retrieval and Spatial Statistics}
To cope with a series of Landsat bands and {search for} metadata, we examined the {image parameters} using regular expression with the command line utilities of Unix as a combination of grep, the Geospatial Data Abstraction Library (GDAL) and echo (Listing \ref{lst:01}). {Here, grep searches for text in a file according to a given pattern of characters, the GDAL filters spatial information in the dataset, and echo outputs the strings using the defined arguments.}

\begin{lstlisting}[language=bash, caption=Unix %MDPI: please check and confirm all the listing format. <-- Answer: yes, we confirm.
shell script by GDAL and echo-grep utilities to inspect the geometry of Landsat TIFs., label={lst:01}]
for f in *tif ; do  echo $f; gdalinfo $f | grep Size; done

\end{lstlisting}

The spatial resolution, corresponding to the pixel size, was 30 m for Bands 1--7, 9 and 10, while it was 15 m for Band 8 (panchromatic) (Figure \ref{fig03}a). 

\begin{figure}[H]

     \begin{subfigure}[b]{0.45\textwidth}
         \centering
         \includegraphics[width=6.3cm]{Fig_03a}
         \caption{Using grep and echo Unix utilities with 'gdalinfo' module to scan the metadata.}
         \label{fig:03a}
     \end{subfigure} \hfill \hspace{2mm}
     \begin{subfigure}[b]{0.45\textwidth}
         \centering
         \includegraphics[width=6.3cm]{Fig_03b}
         \caption{Creating SpatRaster objects in R and checking spatial data: projection, datum and ellipsoid.}
         \label{fig:03b}
     \end{subfigure}
     \vspace{6pt}
        \caption{Inspecting the geometry of Landsat bands{ with }({\bf a}) 30-pixel resolution for all the bands, except for Band 8 (panchromatic) with a 15-pixel resolution, checking image properties in R package 'terra'{ and} ({\bf b}) for single Landsat-8 OLI/TIRS C1 bands (1:11).}
        \label{fig03}
\end{figure}

The spatial information of the Landsat channels (bands) is presented in metadata, as shown in Figure \ref{fig03}b, {and} inspected by R. It indicates the number and location of the spectral measurements, spatial extent, geodetic information (coordinate reference system (CRS), ellipsoid and datum obtained by the global positioning system (GPS) of the spacecraft and the spatial and radiometric resolutions of the image. Each image {corresponding to a given} band of the Landsat scene was read into the R SpatRaster object class, which represents the equalised area of the rectangular cells with regard to the units of the CRS (Figure \ref{fig03}b). 

\subsubsection{Spectral Bands for Colour Composites}

The properties of the image were {evaluated using} a SpatRaster object {by} the R commands presented in Listing \ref{lst:02}. Here, we inspected the {CRS}, the geometry of the image (number of cells, rows and columns) and extent, resolution and origin of the bands. The panchromatic channel (Band 8) had a higher resolution (15 m) compared with the set of the multispectral bands in the visible and thermal bands of the spectra (30 m). To ignore this difference in the matrix of the image, the raster in Band 8 (panchromatic) was resampled to the structure of the same number of rows and columns as in other bands (Here, target Band 1 is in Ultra Blue) by creating {the} RasterBrick multi-layer object {as demonstarted in} Listing \ref{lst:02}. 

\begin{lstlisting}[language=r, caption=R code used for inspecting the information and statistics on the Landsat 8-9 OLI bands., label={lst:02}]
library(terra)
## terra version 1.2.11
# Change working directory
setwd("/Users/polinalemenkova/Documents/R/Image_Processing/LC09_L1TP_197055_20220111_20220122_02_T1")
b2 <- rast('LC09_L1TP_197055_20220111_20220122_02_T1_B2.TIF') # Blue
b3 <- rast('LC09_L1TP_197055_20220111_20220122_02_T1_B3.TIF') # Green
b4 <- rast('LC09_L1TP_197055_20220111_20220122_02_T1_B4.TIF') # Red
b5 <- rast('LC09_L1TP_197055_20220111_20220122_02_T1_B5.TIF') # Near Infrared (NIR)
# Checking metadata by printing the variables
b2
# Checking coordinate reference system (CRS)
crs(b2)
# Number of cells
ncell(b2)
# Number of rows & columns
dim(b2)
# Image information and statistics spatial resolution
res(b2)
# Number of bands
nlyr(b2)
# Comparing bands extent, number of rows and columns, projection, resolution and origin
compareGeom(b5, b4)
# Creating a SpatRaster with multiple layers from the existing SpatRaster (single layer) objects.
s <- c(b5, b4, b3)
# Checking properties of the multi-band image
s
# Creating a list of raster layers to use
filenames <- paste0('LC09_L1TP_197055_20220111_20220122_02_T1_B', 1:11, ".tif")
filenames
# Resampling raster in Band 8 (panchromatic) to the structure of Band 1:
library(raster)
r1 = raster::brick("LC09_L1TP_197055_20220111_20220122_02_T1_B8.tif")
r2 = raster::brick("LC09_L1TP_197055_20220111_20220122_02_T1_B1.tif")
r1 <- resample(r1, r2, method="bilinear")
s <- stack(r1, r2)
# Creating the multi-layer SpatRaster using the filenames
landsat <- rast(filenames)
landsat
\end{lstlisting} 

{Afterwards, }the layers were read in the `landsat' object {class of R,} where they were represented by the bands of {the} Landsat-9 OLI {image.~The} reflection intensity {had} the following wavelengths: Ultra Blue {or} coastal aerosol (Band 1), blue (Band 2), green (Band~3), red (Band 4), near infrared (NIR) (Band 5), Shortwave Infrared (SWIR) 1 (Band 6), Shortwave Infrared (SWIR) 2 (Band 7), panchromatic (Band 8), cirrus (Band 9), Thermal Infrared (TIRS) 1 (Band 10) {and} Thermal Infrared (TIRS) 2 (Band 11). 

The visualisation of single bands was performed using Listing \ref{lst:03} and is presented as a subplot of grayscale scenes in Figure \ref{fig04}. {The }visual inspection of the satellite image product built from the monochrome image data (Figure \ref{fig04}) was useful for extracting general information on {the} large-scale features,{ }such as river systems, lakes, mountains, fracture patterns, large geological lineaments {or the extent of the city} with notable objects{ such as }highways. These objects are visible and can be identified by their shape and size and interpreted using spatial properties related to the{ }pixel brightness{ values}. 

\begin{lstlisting}[language=r, caption=R code for plotting the 11 original individual layers (raw bands) of the Landsat-9 multi-spectral images as grayscale bitmap images., label={lst:03}]
library(terra)
## terra version 1.2.11
# Change working directory
setwd("/Users/polinalemenkova/Documents/R/Image_Processing/LC09_L1TP_197055_20220111_20220122_02_T1")
b1 <- rast('LC09_L1TP_197055_20220111_20220122_02_T1_B1.TIF') # Ultra Blue
b2 <- rast('LC09_L1TP_197055_20220111_20220122_02_T1_B2.TIF') # Blue
b3 <- rast('LC09_L1TP_197055_20220111_20220122_02_T1_B3.TIF') # Green
b4 <- rast('LC09_L1TP_197055_20220111_20220122_02_T1_B4.TIF') # Red
b5 <- rast('LC09_L1TP_197055_20220111_20220122_02_T1_B5.TIF') # NIR
b6 <- rast('LC09_L1TP_197055_20220111_20220122_02_T1_B6.TIF') # Green
b7 <- rast('LC09_L1TP_197055_20220111_20220122_02_T1_B7.TIF') # Red
b8 <- rast('LC09_L1TP_197055_20220111_20220122_02_T1_B8.TIF') # Panchrom.
b9 <- rast('LC09_L1TP_197055_20220111_20220122_02_T1_B9.TIF') # Cirrus
b10 <- rast('LC09_L1TP_197055_20220111_20220122_02_T1_B10.TIF') # TIRS 1
b11 <- rast('LC09_L1TP_197055_20220111_20220122_02_T1_B11.TIF') # TIRS 2.
# Plotting natural color composite from bands B4, B3 and B2.
landsatRGB <- c(b4, b3, b2)
plotRGB(landsatRGB, stretch = "lin")
# Plotting 11 original individual layers as raw bands of the Landsat-9 OLI.
par(mfrow = c(4, 3))
plot(b1, main = "Ultra Blue (Band 1)", col = gray(0:100 / 100), legend=FALSE)
plot(b2, main = "Blue (Band 2)", col = gray(0:100 / 100), legend=FALSE)
plot(b3, main = "Green (Band 3)", col = gray(0:100 / 100), legend=FALSE)
plot(b4, main = "Red (Band 4)", col = gray(0:100 / 100), legend=FALSE)
plot(b5, main = "NIR (Band 5)", col = gray(0:100 / 100), legend=FALSE)
plot(b6, main = "SWIR (Band 6)", col = gray(0:100 / 100), legend=FALSE)
plot(b7, main = "NIR (Band 7)", col = gray(0:100 / 100), legend=FALSE)
plot(b8, main = "Panchrom. (Band 8)", col = gray(0:100 / 100), legend=FALSE)
plot(b9, main = "Cirrus (Band 9)", col = gray(0:100 / 100), legend=FALSE)
plot(b10, main = "TIRS 1 (Band 10)", col = gray(0:100 / 100), legend=FALSE)
plot(b11, main = "TIRS 2 (Band 11)", col = gray(0:100 / 100), legend=TRUE)
landsatRGB <- c(b4, b3, b2)
plotRGB(landsatRGB, stretch = "lin")
\end{lstlisting} 

There was a difference in {nuance} between the grey corresponds and the numerical values and range of pixel values in diverse channels of Landsat.%Check meaning retained.
~{This} is {explained by} the spectral reflectance that differs in various bands for objects on the Earth's surface. Due to the individual reflectance of{ }solar radiation by {these} features, the values differed accordingly. {Such effects from spectral reflectance }enabled {us} to use combinations of various bands {to} better represent{ the }diverse phenomena and objects on Earth{. }For {instance}, the particular applications of Landsat bands are for geology (SWIR (Band 7), NIR (Band 5) and red (Band~4)) \cite{VOLESKY2003235}, mineralogical identification of soil and rock (SWIR (Band 7), red (Band 4) and green (Band 2)) and identifying {the} difference between soil and vegetation (green in Band 3) \cite{Richards,CampbellWynne}. Aside from that, the reflectance data from the selected bands of the Landsat images were used in the calculation of{ }four vegetation indices for the study area using the interactions of pixels in the raster.

Subsetting large files, such as multispectral Landsat {images}, was {performed} using {the} SpatRaster object and 'subset' function, {which enabled} processing{ }the necessary bands{ of the Landsat image}. In this respect, it enabled ignoring{ }the cirrus, TIRS and other{ }bands{ that were not relevant  for a given task }and to leave only{ }the visible spectra. In this way, a subset of {the} target layers was selected from the SpatRaster object, while{ }other bands were removed from the dataset{ }using Listing \ref{lst:04}.
\vspace{-4pt}
\begin{figure}[H]
\makebox[\textwidth][r]{\includegraphics[width=1.3\textwidth]{Fig_04.jpg}}%
	\caption{\hl{Original individual} %MDPI: Please use commas to separate thousands for numbers with five or more digits (not four digits) in the picture, e.g., "10000" should be "10,000".
%MDPI: Please change the terms into scientific notations in the figure,  e.g., "$8 \times 10^{3}$", not "8E3".
 11 layers (raw bands) of the multi-spectral image Landsat-9 OLI/TIRS C1 `LC09\_L1TP\_197055\_20220111\_20220122\_02\_T1' in greyscale mode: ultra blue (Band 1), blue (Band 2), green (Band 3), red (Band 4), near infrared (NIR) (Band 5), Shortwave Infrared (SWIR) 1 (Band 6), Shortwave Infrared (SWIR) 2 (Band 7), panchromatic (Band 8), cirrus (Band 9), Thermal Infrared (TIRS) 1 (Band 10) and Thermal Infrared (TIRS) 2 (Band 11), plotted with R.}
	\label{fig04}
\end{figure} 
\newpage
\begin{lstlisting}[language=r, caption=R code for subsetting the large files using SpatRaster object., label={lst:04}]
library(terra)
setwd("/Users/polinalemenkova/Documents/R/Image_Processing/LC09_L1TP_197055_20220111_20220122_02_T1")
filenames <- paste0('LC09_L1TP_197055_20220111_20220122_02_T1_B', 1:11, ".tif")
filenames
landsat <- rast(filenames)
landsat
# Removing the unnecessary four bands: Panchromatic, Cirrus, TIRS1, TIRS2.
landsatsub <- subset(landsat, 1:7)
names(landsatsub)
# Checking Nnumber of bands in the original and new data
nlyr(landsatsub)
# Renaming the bands
names(landsatsub) <- c("ultra-blue", "blue", "green", "red", "NIR", "SWIR1", "SWIR2")
names(landsatsub)
\end{lstlisting} 

To more thoroughly investigate the relationship between the spectral responses of various land cover types in different wavelength diapasons, we examined the pairwise correlation between the selected bands{ of the Landsat image}. To this end, we compared the bands in pairwise combinations (Band 1 {against} Band 2 and Band 4 {against} Band 5) in the selected parts of the spectrum{. The }images were plotted pixel by pixel against those from the corresponding bands{ using }Listing \ref{lst:05}. The resulting plot showed a narrow correlation ellipse and a histogram representing a strong positive correlation between the two bands of green and blue as well as red to NIR (Figure \ref{fig05}).

\begin{lstlisting}[language=r, caption=R code for inspecting the correlation between the selected Landsat bands., label={lst:05}]
# Relation between bands
library(terra)
pairs(landsatsub[[1:2]], hist=TRUE, cor=TRUE, use="pairwise.complete.obs", maxcells=100000, main = "Correlation of Band 1 (Ultra-blue) against Blue (Band 2)")
pairs(landsatsub[[4:5]], hist=TRUE, cor=TRUE, maxcells=100000, main = "Correlation of Band 4 (Red) against NIR (Band 5)")
\end{lstlisting} 

The cells with low spectral reflectance in a target band {generally} had low reflectance in{ }the {remaining} bands, which was caused by the albedo of the Earth's surface. However, {a visible} correlation between Bands 1 and 2 showed more similarity between these bands, {as} all the pixels {were located} almost {along} the same line{ (Figure} \ref{fig05}a){. In contrast, the }NIR and red bands demonstrated more differences in spectral reflectance for the selected objects, which resulted in {a} scattered cloud of points representing the pixels {of the image (Figure} \ref{fig05}b). The brightness of the pixels in each layer indicates the amount of  incident solar radiation, which is reflected by the Earth's surface in a given wavelength diapason{. As a result, this causes a}  similarity {or} difference in the spectral bands of the Landsat scene. 

Since there was not much information in the monochrome bands {of the Landsat OLI/TIRS image}, they were combined into {colour} composites to create more {representative images}. Plotting{ of the }colour composites was based on the combinations of { }the three{ }bands used as elements of the matrix scene. The three bands assigned to each of the three primary colours (red, green and blue (RGB)) were displayed with images {using} Listing \ref{lst:06}.
\begin{figure}[H]

     \begin{subfigure}[b]{0.45\textwidth}
      
         \includegraphics[width=6.5cm]{Fig_05a.jpg}
         \caption{   \centering}
     \end{subfigure} \hfill \hspace{2mm}
     \begin{subfigure}[b]{0.45\textwidth}
   
         \includegraphics[width=6.5cm]{Fig_05b.jpg}
         \caption{   \centering}
     \end{subfigure}
        \caption{Correlation %MDPI: Please use commas to separate thousands for numbers with five or more digits (not four digits) in the picture, e.g., "10000" should be "10,000". <--  Answer: corrected (Figures 5a and 5b are replotted).
between the selected Landsat-8 OLI/TIRS C1 bands in a multispectral image \emph{LC09\_L1TP\_197055\_20220111\_20220122\_02\_T1\_B}. ({\bf a}) Band 1 against Band 2. %MDPI: We moved the subfigure explanations into the figure caption. Please confirm. <-- Answer: yes, we confirm.
 ({\bf b}) Band 4 against Band~5.}
        \label{fig05}
\end{figure}
\vspace{-12pt}


\begin{lstlisting}[language=r, caption=R code used for plotting the color composites for bands of the satellite image., label={lst:06}]
library(raster)
blue = raster("LC09_L1TP_197055_20220111_20220122_02_T1_B2.TIF")
green = raster("LC09_L1TP_197055_20220111_20220122_02_T1_B3.TIF")
red = raster("LC09_L1TP_197055_20220111_20220122_02_T1_B4.TIF")
rgb = stack(red, green, blue)
plotRGB(rgb, scale=65535)
plotRGB(rgb, stretch="lin")
plotRGB(rgb, stretch="hist")
# False color composites
# Color Infrared: colors near-infrared (band 5) as red, red (band 4) as green, and green (band 3) as blue
green = raster("LC09_L1TP_197055_20220111_20220122_02_T1_B3.TIF")
red = raster("LC09_L1TP_197055_20220111_20220122_02_T1_B4.TIF")
near.infrared = raster("LC09_L1TP_197055_20220111_20220122_02_T1_B5.TIF")
rgb = stack(near.infrared, red, green)
plotRGB(rgb, scale=65535, stretch="lin")
# Landsat 9 OLI/TIRS colour composite image, RGB=753
green = raster("LC09_L1TP_197055_20220111_20220122_02_T1_B7.TIF")
red = raster("LC09_L1TP_197055_20220111_20220122_02_T1_B5.TIF")
near.infrared = raster("LC09_L1TP_197055_20220111_20220122_02_T1_B3.TIF")
rgb = stack(near.infrared, red, green)
plotRGB(rgb, scale=65535, stretch="hist")
# Landsat 9 OLI/TIRS colour composite image, RGB=765.
red = raster("LC09_L1TP_197055_20220111_20220122_02_T1_B7.TIF")
green = raster("LC09_L1TP_197055_20220111_20220122_02_T1_B6.TIF")
blue = raster("LC09_L1TP_197055_20220111_20220122_02_T1_B5.TIF")
rgb = stack(red, green, blue)
plotRGB(rgb, scale=65535, stretch="hist")
# Landsat 9 OLI/TIRS colour composite image, RGB=742.
red = raster("LC09_L1TP_197055_20220111_20220122_02_T1_B7.TIF")
green = raster("LC09_L1TP_197055_20220111_20220122_02_T1_B4.TIF")
blue = raster("LC09_L1TP_197055_20220111_20220122_02_T1_B2.TIF")
rgb = stack(red, green, blue)
plotRGB(rgb, scale=65535, stretch="hist")
# Landsat 9 OLI/TIRS colour composite image, RGB=573.
red = raster("LC09_L1TP_197055_20220111_20220122_02_T1_B5.TIF")
green = raster("LC09_L1TP_197055_20220111_20220122_02_T1_B7.TIF")
blue = raster("LC09_L1TP_197055_20220111_20220122_02_T1_B3.TIF")
rgb = stack(red, green, blue)
plotRGB(rgb, scale=65535, stretch="hist")
# Landsat 9 OLI/TIRS colour composite image, Color Infrared: RGB=543
green = raster("LC09_L1TP_197055_20220111_20220122_02_T1_B5.TIF")
red = raster("LC09_L1TP_197055_20220111_20220122_02_T1_B4.TIF")
near.infrared = raster("LC09_L1TP_197055_20220111_20220122_02_T1_B3.TIF")
rgb = stack(near.infrared, red, green)
plotRGB(rgb, scale=65535, stretch="hist")
# Landsat 9 OLI/TIRS colour composite image, Color Infrared: RGB=654
green = raster("LC09_L1TP_197055_20220111_20220122_02_T1_B6.TIF")
red = raster("LC09_L1TP_197055_20220111_20220122_02_T1_B5.TIF")
near.infrared = raster("LC09_L1TP_197055_20220111_20220122_02_T1_B4.TIF")
rgb = stack(near.infrared, red, green)
plotRGB(rgb, scale=65535, stretch="hist")
\end{lstlisting}

The vegetation indices were computed based on the arithmetically transformed combination of Landsat OLI/TIRS bands. This resulted from the existing formulae for band computations using mathematical operations {with bands.~The }map algebra included{ }addition, subtraction and division of the brightness{ in cells }of the NIR and red bands of the Landsat image. The differences in the bands highlight the regions of the enhanced and more healthy vegetation in the target area of the image, which {presents information} for environmental monitoring{.} The resultant image shows{ }brightness values of vegetation modified by colour palettes made using arithmetic operations on the composite bands. 

The most widely used and well-known vegetation index is the Normalised Difference Vegetation Index (NDVI), which applies the ratio of the difference in values between {the} red and NIR bands and their sum. Similar to the{ }NDVI, {the} Enhanced Vegetation Index (EVI) quantifies vegetation greenness with added corrections for the atmospheric conditions and canopy background noise. Aside from that, the EVI is more sensitive in{ }forest areas with{ }more dense canopies and vegetation coverage,{ }which is useful for{ }the forest regions of Côte d'Ivoire. The Soil-Adjusted Vegetation Index (SAVI) is built up by using the NDVI and {correcting for} the effects of soil brightness in the regions where vegetative cover is low{, namely }deforested areas{ }and the fragmented savannahs of Côte d'Ivoire. The Atmospherically Resistant Vegetation Index 2 (ARVI2) is the most applicable index in regions with high atmospheric aerosol contents. Examples of other environmental indices include{ }the Normalized Difference Water Index (NDWI), modified Normalized Difference Water Index (MNDWI) and Water Ratio Index (WRI) \cite{s22186827}.



Computing the vegetation indices is an important environmental tool which is based on information from {the DN values of pixels}, indicating the spectral reflectance of the red and NIR bands as measured by the sensor. It is commonly computed with the spatial statistics and algebra of the bands using the cell values {of the image} as {shown} in Listing \ref{lst:07}. The reason for selecting the NIR and red bands for computing the vegetation indices is vegetation appearing brighter in these bands, since it reflects more energy in the NIR and red bands compared {with} other wavelengths. Contrarily, water absorbs most of the energy in the NIR band and is {therefore} represented by black to very dark brown colours.~Through a series of interactions among the{ }elements{ of an image}, they are allocated to the classes based on the weighted values of their relative proportions.
\begin{lstlisting}[language=r, caption=R code used for computing vegetation indices from bands of the satellite image., label={lst:07}]
# Computing vegetation indices
library(terra)
library(RColorBrewer)
setwd("/Users/polinalemenkova/Documents/R/Image_Processing/LC09_L1TP_197055_20220111_20220122_02_T1")
# 1. Normalized Difference Vegetation Index (NDVI) = (NIR - R) / (NIR + R) NDVI = (Band 5 - Band 4) / (Band 5 + Band 4).
vi <- function(img, k, i) {
  bk <- img[[k]]
  bi <- img[[i]]
  vi <- (bk - bi) / (bk + bi)
  return(vi)
}
# For Landsat NIR = 5, red = 4.
ndvi <- vi(landsat, 5, 4)
plot(ndvi, col=rev(terrain.colors(20)), font.main = 1, main = "NDVI for Landsat-9 OLI/TIRS C1 image LC09_L1TP_197055_20220111_20220122_02_T1: Yamassoukro, central C\^ote d'Ivoire", cex.main=0.9)
# Plotting histogram of the NDVI
hist(ndvi, font.main = 1, main = "NDVI values for Landsat-9 OLI/TIRS C1 image \n LC09_L1TP_197055_20220111_20220122_02_T1, Yamassoukro, central C\^ote d'Ivoire", xlab = "NDVI", ylab= "Frequency",
    col = "darkolivegreen1", xlim = c(-0.5, 1),  breaks = 30, xaxt = "n")
axis(side=1, at = seq(-0.6, 1, 0.1), labels = seq(-0.6, 1, 0.1))
#
# 2. Soil Adjusted Vegetation Index (SAVI) corrects NDVI for the influence of soil brightness in areas where vegetative cover is low
# SAVI = ((Band 5 - Band 4) / (Band 5 + Band 4 + 0.5)) * (1.5).
savi_fun = function(nir, red){
  ((nir - red) / (nir + red + 0.5)) * 1.5
}
savi = lapp(landsat[[c(4, 3)]],
            fun = savi_fun)
#plot(savi, col = topo.colors(30), font.main = 2, main = "Soil Adjusted Vegetation Index (SAVI) for Landsat-9 OLI/TIRS C1 image \nLC09_L1TP_197055_20220111_20220122_02_T1: Yamassoukro, central C\^ote d'Ivoire", cex.main=1.0)
plot(savi, col = brewer.pal(11, "Accent"), font.main = 2, main = "Soil Adjusted Vegetation Index (SAVI) for Landsat-9 OLI/TIRS C1 image \nLC09_L1TP_197055_20220111_20220122_02_T1: Yamassoukro, central C\^ote d'Ivoire", cex.main=1.0)
# Computing histogram for SAVI
hist(savi, font.main = 1, main = "SAVI values for Landsat-9 OLI/TIRS C1 image \n LC09_L1TP_197055_20220111_20220122_02_T1, Yamassoukro, central C\^ote d'Ivoire", xlab = "SAVI", ylab= "Frequency",
    col = "darkgoldenrod1", xlim = c(-0.4, 0.4),  breaks = 30, xaxt = "n")
axis(side=1, at = seq(-0.4, 0.4, 0.1), labels = seq(-0.4, 0.4, 0.1))
#
# 3. Enhanced Vegetation Index 2
# EVI2 = 2.4 * (Band 5 - Band 4) / (Band 5 + Band 4 + 1).
evi2_fun = function(nir, red){
  2.4 * ((nir - red) / (nir + red + 1))
}
evi2 = lapp(landsat[[c(4, 3)]],
            fun = evi2_fun)
plot(evi2, col = brewer.pal(9, "Set1"), font.main = 2, main = "Enhanced Vegetation Index 2 (EVI2) for Landsat-9 OLI/TIRS C1 image \nLC09_L1TP_197055_20220111_20220122_02_T1: Yamassoukro, central C\^ote d'Ivoire", cex.main=0.9)
# Computing histogram for EVI2
hist(evi2, font.main = 1, main = "EVI2 values for Landsat-9 OLI/TIRS C1 image \n LC09_L1TP_197055_20220111_20220122_02_T1, Yamassoukro, central C\^ote d'Ivoire", xlab = "SAVI", ylab= "Frequency",
    col = "darkorchid1", xlim = c(-0.4, 0.4),  breaks = 30, xaxt = "n")
axis(side=1, at = seq(-0.4, 0.4, 0.1), labels = seq(-0.4, 0.4, 0.1))
#
# 4. Atmospherically Resistant Vegetation Index 2 (ARVI2)
# ARVI2 = -0.18 + 1.17 ((nir - red) / (nir + red))
arvi_fun = function(nir, red){
  ((nir - red) / (nir + red)) * 1.17 - 0.18
}
arvi = lapp(landsat[[c(4, 3)]],
            fun = arvi_fun)
plot(arvi, col = brewer.pal(12, "Paired"), font.main = 1, main = "Atmospherically Resistant Vegetation Index 2 (ARVI2) for Landsat-9 OLI/TIRS C1 \nimage LC09_L1TP_197055_20220111_20220122_02_T1: Yamassoukro, C\^ote d'Ivoire", cex.main=0.9)
# Computing histogram for ARVI2
hist(arvi, font.main = 1, main = "ARVI2 values for Landsat-9 OLI/TIRS C1 image \n LC09_L1TP_197055_20220111_20220122_02_T1, Yamassoukro, C\^ote d'Ivoire", xlab = "ARVI2", ylab= "Frequency",
    col = "cyan1", xlim = c(-0.4, 0.1),  breaks = 30, xaxt = "n")
axis(side=1, at = seq(-0.4, 0.1, 0.1), labels = seq(-0.4, 0.1, 0.1))
\end{lstlisting}

 
\subsection{Terrain Analysis in Python}
To model the relationship between the topographic variables and appearance, we applied Python algorithms for processing SRTM data by the 3D approach. {Two different types of local features were computed for the SRTM image, namely the DTM and hillshade, with artificial light imitating the Sun's illumination.} The 3D visualisation of the terrain surface is {changed by varied} azimuth and altitude of sun{. }Technically, the terrain analysis has been performed using integration of {several libraries of} Python{: }EarthPy \cite{wasser2019EarthPyAP}{ for modelling}, Matplotlib for data visualisation \cite{Hunter2007MatplotlibA2}, Rasterio for processing gridded raster datasets \cite{gillies_2019}{. The }auxiliary packages {include }NumPy \cite{harris2020array} {which} operates with arrays{ }of raster cells{ presenting the elements of the matrix of image, arranged in regular rows and columns,} and Plotly \cite{plotly},{ which} is used as {an} additional graphical and plotting library.

The geospatial{ }TIFF was visualised for modelling the DEM as GeoTIFF using the Rasterio library, which {imported }the dataset as a raster array of cells, read the GeoTIFF format and stored a gridded raster{ as} a digital terrain model{ }(DTM) {or} satellite images (Listing \ref{lst:08}). {The} elevation data were {imported by Rasterio} as height values that were contained in the cells of the corresponding raster.~{The structure of the data presented a NumPy array in the form of a matrix with the DNs of pixels}. The terrain was{ }plotted as an array of cells representing the land surface according to the elevation values, as shown in Listing \ref{lst:08}. {Plotting of the bands  was performed using the ep.plot\_bands() function using the spatial dimensions of the raster file.}

\begin{lstlisting}[language=python, caption=Python code (1) used for plotting DEM in Figure \ref{fig02}., label={lst:08}]
python3.10
# Set home directory
import os
os.chdir('/Users/polinalemenkova/Documents/Python/Image_Processing')
os. getcwd()
# Set data source
dtm = "n07_w006_3arc_v2.tif"
# Load libraries
from glob import glob
import earthpy as et
import earthpy.spatial as es
import earthpy.plot as ep
import rasterio as rio
import matplotlib.pyplot as plt
import numpy as np
from matplotlib.colors import ListedColormap
import plotly.graph_objects as go
np.seterr(divide='ignore', invalid='ignore')
# Open DEM by Rasterio 
with rio.open(dtm) as src:
    elevation = src.read(1)
# Plot the data
fig, ax = plt.subplots(figsize=(10, 10))
ep.plot_bands(
    elevation,
    ax=ax,
    cmap="terrain",
    title="Ditigal Terrain Model (DTM) of Kossou Lake, C\^ote d'Ivoire. \n Data source: SRTM3 (3 arc-second resolution, 90 m), NASA/NGA. \nEntity ID: SRTM3N07W006V2. Aquisition date: 2000-02-11. cpt: 'terrain'",
    figsize=(10, 6),
)
plt.show()
# Save figure with 300 dpi resolution
plt.savefig('Figure_2_DEM_Kossou_1609.jpg', dpi=300)
\end{lstlisting}

The hillshade was calculated using the altitude of the Sun as a light source at the zenith and azimuth of the illumination source (Listing \ref{lst:09}). {The functions to compute the variations between the exemplar locations of artificial illumination of the terrain resulted in the visualised details of the terrain in the Kossou Lake surroundings. }The slope and aspect of the terrain were used as auxiliary variables for the terrain modelling. The hillshade modelling was inherently related to static cartographic representations of the Earth's topography, and thus hillshade was used for the geomorphometric representation of the terrain of Côte d'Ivoire. 

\begin{lstlisting}[language=python, caption=Python code (2) used for plotting hillshade in Figure \ref{fig03}., label={lst:09}]
python3.10
# Set home directory
import os
os.chdir('/Users/polinalemenkova/Documents/Python/Image_Processing')
os. getcwd()
# Set data source
dtm = "n07_w006_3arc_v2.tif"
# Load libraries
from glob import glob
import earthpy as et
import earthpy.spatial as es
import earthpy.plot as ep
import rasterio as rio
import matplotlib.pyplot as plt
import numpy as np
from matplotlib.colors import ListedColormap
import plotly.graph_objects as go
np.seterr(divide='ignore', invalid='ignore')
# Generate hillshade from DTM
hillshade = es.hillshade(elevation)
# Plot the data
fig, ax = plt.subplots(figsize=(10, 10))
ep.plot_bands(
    hillshade,
    ax=ax,
    cbar=True,
    cmap="cividis",
    title="Hillshade from DTM of Kossou Lake, C\^ote d'Ivoire. \n Data source: SRTM3 (3 arc-second resolution, 90 m), NASA/NGA. \nEntity ID: SRTM3N07W006V2. Aquisition date: 2000-02-11. cmap=cividis",
    figsize=(10, 10),
)
plt.show()
plt.savefig('Figure_3_Hillshade_Kossou_1609.png', dpi=300)
\end{lstlisting}

Hillshade represents the terrain with a hypothetical illumination value in the monochrome scale (0--255) with cells on the raster scene showing the terrain surface with a given specified light source. The illumination was changed to 10\degree and 230\degree~(in radians) to better represent a geospatial model of the topographic relief of Côte d'Ivoire (Listing \ref{lst:10}). {Assume that the orientation differences between the angle and azimuth of artificial illumination in modelling the topography of Kossou Lake are independent and identically distributed according to the time and seasonal parameters. To visualise the characteristics of the relief in the topographic model based on the DEM, raster data were sampled to estimate the variations in distribution of the parameters. In our case, we selected azimuths of 10\degree and 230\degree} (\hl{Figure}%MDPI: Please cite the figure in the text and ensure the first citation of each figure appears in numerical order. Figure 13 is after Figure 5, please confirm and revise
~\ref{fig13}c,d). We observed that the effects from the contrasting variations in relief approximately followed the changes in the Sun's azimuth and angle altitude.

\begin{lstlisting}[language=python, caption=Python code used for plotting the hillshade with changed azimuth and Sun angle values for Figures \ref{fig04} and \ref{fig05}., label={lst:10}]
python3.10
# Set home directory
import os
os.chdir('/Users/polinalemenkova/Documents/Python/Image_Processing')
os. getcwd()
# Set data source
dtm = "n07_w006_3arc_v2.tif"
# Load libraries
from glob import glob
import earthpy as et
import earthpy.spatial as es
import earthpy.plot as ep
import rasterio as rio
import matplotlib.pyplot as plt
import numpy as np
from matplotlib.colors import ListedColormap
import plotly.graph_objects as go
np.seterr(divide='ignore', invalid='ignore')
# Open DEM by Rasterio which reads GeoTIFF format and stores gridded raster datasets of digital terrain models (DTMs) and satellite images
with rio.open(dtm) as src:
    elevation = src.read(1)
# Changing the azimuth of the sun of the hillshade layer to 210 degrees
hillshade_azimuth_230 = es.hillshade(elevation, azimuth=230)
# Plotting the hillshade layer with the modified azimuth
fig, ax = plt.subplots(figsize=(10, 10))
ep.plot_bands(
    hillshade_azimuth_230,
    ax=ax,
    cbar=True,
    cmap="plasma",
    title="Hillshade from DTM of Kossou Lake, C\^ote d'Ivoire. Azimuth of sun is adjusted to 230 \N{DEGREE SIGN}",
    figsize=(10, 10),
)
plt.show()
plt.savefig('Figure_4_Azimuth_hillshade_Kossou.jpg', dpi=300)
# Changing the angle altitude of sun
hillshade_angle_10 = es.hillshade(elevation, altitude=10)
# Plotting the hillshade layer with the modified angle altitude
fig, ax = plt.subplots(figsize=(10, 10))
ep.plot_bands(
    hillshade_angle_10,
    ax=ax,
    cbar=True,
    cmap="magma",
    title="Hillshade from DTM of Kossou Lake, C\^ote d'Ivoire. Angle altitude is defined as 10 \N{DEGREE SIGN}",
    figsize=(10, 10),
)
plt.show()
plt.savefig('Figure_5_Angle_hillshade_Kossou.jpg', dpi=300)
\end{lstlisting}

{We }assumed the data to be of a high resolution, {which may have presented} our current limitations{. }However, {the precision }may affect the resulting topographic model through coarser{ }resolution data (i.e., lesser meters in one pixel will give a rough representation of the Earth's surface and miss smaller landforms in the landscape). Regarding this, it is strongly recommended to apply {very }high-resolution data for a precise representation of the Earth. Another point of concern {consists of selecting the} correct view shed and azimuth for modelling the aspect, slope and elevation, which is {especially} true for{ }mountainous regions and rugged reliefs.{ }Although the Python-based method of terrain modelling can potentially handle materials with coarser or unknown resolutions for the original data, the results may {hide the} selected structures on the Earth's surface. The same {issue} concerns the applications of R for image processing, {where the} precision {of the output data corresponds to } to the original resolution of the input {data}. In our future works, we are very interested in extending our {study} to handle higher-resolution imagery, such as Sentinel 2A at a 10 m resolution or Satellite pour l'Observation de la Terre (SPOT) with up to a 1.5 m spatial resolution in the panchromatic and multispectral channels.%Check meaning retained.

{In the presented methodology, we} applied Python and R algorithms which {process and model} the terrain characteristics of the Earth {using the spatial parameters of the remote sensing data}. By exploiting the {spectral reflectance} characteristics of the satellite bands, we {used information} from {the Landsat channels from} the {corresponding} cells which {represent} objects located on the Earth.~Interpretation of these values is based on their spectral reflectance in various wavebands, which {provides} information {on} the characteristics of typical land cover types. This was possible through the analysis of the different brightness of each land cover type (e.g., forests, rivers, lakes and other water bodies {or} urban spaces). In turn, the geomorphometric models formed using Python scripts enabled predicting the distribution of various soil properties and associated vegetation types which were related to the specific slope steepness and regional topographic features. {Examples of} the application of terrain models {include} hazard risk assessment in applied geology{ for }erosion, landslides and avalanches{ }in mountainous regions. The major advantage of using Python and R scripts for environmental modelling is that it reduces the time of data processing and increases the precision of models through automation of data processing, {because} scripts {are repeatable} for other regions and {can be }applied in similar studies.

%%%%%%%%%%%%%%%%%%%%%%%%%%%%%%%%%%%%%%%%%%
%\newpage
\section{Results and Discussion}
%--------
\subsection{Color Composites}
Colour image maps were created by different combinations of the recorded Landsat OLI/TIRS bands. In the sequence from top to bottom, they are displayed as natural and false composites,  respectively{ }(Figure \ref{fig06}).~The use of natural (or true) colour composites was widely applied for general environmental monitoring, since it exposed{ }characteristics regarding the landscape and land cover types compared to the monochrome images.{ }Thus, band combinations 4-3-2 (red-green-blue) and 2-3-4 (blue-green-red) represented natural colour composites well. Since its appearance is similar to the normal aerial photograph, this composite is advantageous for mapping vegetation and agricultural coverage and discriminating the vegetation from bare soil and growing crops (Figure \ref{fig06}a,b). 

In Figure \ref{fig06}c, Bands 3,4 and 5 were taken in the inverse way, where Band 5 (NIR) represented red,  Band 4 (red) was green and  Band 3 (green) was blue. The improvement of the image was performed using linear and histogram stretching, which consisted of the adjustment of the contrast in the colour tones of the images. In this band composite, the broadleaf healthy vegetation had a deeper reddish hue, while lighter hues were related to the grasslands, savannahs{ }or sparsely vegetated areas. Soil was represented by various hues, depending on the soil type, from dark to light browns. The populated urban areas were shown in light grey to white spots. 

\textls[-15]{The false colour composite{ was }represented by the mix of bands 3, 4 and 5, where the NIR, red, and green bands were {used. This} printed the optical data, with vegetation dominating through the bright red hues due to the chlorophyll reflection, while agricultural lands were depicted in cyan. The bare soils in this combination appeared{ }as shades of white to light grey and blueish or green in the agricultural regions, depending on the crop type. The darker shades of each colour {generally} indicated moister soil and cities with hues of steel grey (Figure \ref{fig06}c). The combination of Bands 7,6 and 5 showed the mix of SWIR 2, SWIR 1 and NIR. Here, the crop lands were coloured yellow, while vegetation areas were shown as bright ultramarine blue {colour} (Figure \ref{fig06}d). }



Figure \ref{fig07} shows different cases of false colour composites mixed using NIR, SWIR2 and green (a), SWIR 2, red and blue (b), NIR, red and green (c) and SWIR 1, red and green (d). Since NIR light is reflected from the surfaces with higher levels of vegetation{ }but{ }absorbed by water, such band combinations are useful for vegetation and agriculture applications (e.g.{, }mapping forests and crops).
\begin{figure}[H]
\makebox[0.90\linewidth][c]{%
	\begin{subfigure}[b]{.6\textwidth}
	
			\includegraphics[width=0.9\textwidth]{Fig_06a}
		\caption{\centering Bands 2,3 and 4 with stretch improvement: 'linear'}
	\end{subfigure}%
	\begin{subfigure}[b]{.6\textwidth}
	
		\includegraphics[width=0.9\textwidth]{Fig_06b}
		\caption{\centering Bands 2, 3 and 4 with stretch improvement: 'histogram'}
	\end{subfigure}%
}\\
\vfill \vspace{1mm}
\makebox[0.90\linewidth][c]{%
	\begin{subfigure}[b]{.6\textwidth}
	
			\includegraphics[width=0.9\textwidth]{Fig_06c}
		\caption{\centering Bands 3, 4 and 5}
	\end{subfigure}%
	\begin{subfigure}[b]{.6\textwidth}

		\includegraphics[width=0.9\textwidth]{Fig_06d}
		\caption{\centering Bands 7, 6 and 5}
	\end{subfigure}%
}
\caption{Natural color composites ({\bf a},{\bf b}). False color composites ({\bf c},{\bf d}). Landsat 9 OLI image of Kossou Lake region and Yamoussoukro, Côte d'Ivoire. Mapping: RStudio. Source: authors.}\label{fig06}
\end{figure}



\begin{figure}[H]
\makebox[0.90\linewidth][c]{%
	\begin{subfigure}[b]{.6\textwidth}
		
			\includegraphics[width=0.89\textwidth]{Fig_07a}
		\caption{\centering }
	\end{subfigure}%
	\begin{subfigure}[b]{.6\textwidth}
		
			\includegraphics[width=0.89\textwidth]{Fig_07b}
		\caption{\centering }
	\end{subfigure}%
}\\
 \vfill \vspace{1mm}
\makebox[0.90\linewidth][c]{%
	\begin{subfigure}[b]{.6\textwidth}
	
			\includegraphics[width=0.89\textwidth]{Fig_07c}
		\caption{	\centering }
	\end{subfigure}%
	\begin{subfigure}[b]{.6\textwidth}
	
			\includegraphics[width=0.89\textwidth]{Fig_07d}
		\caption{\centering }
	\end{subfigure}%
}
\caption{False color composite visualisation of Landsat 8 OLI image of Kossou Lake surroundings in Côte d'Ivoire. Mapping: RStudio. Source: authors. ({\bf a}) Bands 5, 7 and 3. %MDPI: We moved the subfigure explanations into the figure caption. Please confirm. <-- Answer: yes, we confirm.
 ({\bf b}) Bands 7, 4 and 2. ({\bf c}) Bands 5, 4 and 3 (color infrared). ({\bf d}) Bands 6, 5 and 4.}\label{fig07}
\end{figure}

\subsection{Computing Vegetation Indices}
\subsubsection{Normalized Difference Vegetation Index (NDVI)}

The NDVI aims at quantifying the concentrations of the green leaf canopy. It indicates healthy, vigourous vegetation by using the arithmetic of  the red and NIR bands as shown in Equation (\ref{eq:1})~\cite{1974NASSP}:

\begin{equation}\label{eq:1}
NDVI=\dfrac{(NIR-Red)}{(NIR+Red)}
\end{equation}

Originally tested in 1974 \cite{1974NASSP}, the NDVI received wide recognition in image processing for environmental modelling{. Nowadays, the NDVI} is the most famous and most used vegetation index. The selection of the red and NIR bands in the NDVI is explained by the specific certain abilities of plants to selectively absorb, reflect or transmit the frequencies in the wavebands; plants absorb red and strongly reflect NIR waves. Practically, the NDVI enables indicating the green vegetation density covering the terrain surface and identifying where plants are thriving and where they are under stress (i.e., due to a lack of water).

{Due to} this fundamental property, the NDVI was used to evaluate vegetation health and assess the vigour of plants in agricultural regions. The general NDVI lies in the diapason from $-$1 to +1, but the majority of the values are restricted by a range from $-$0.1 to 0.5 (Figure \ref{fig08}). Negative values for the NDVI correspond to water or very moist soil areas, as shown by the beige-to-white colours in Figure \ref{fig08}, while higher NDVI values signify a dense vegetation canopy, which is depicted by the bright green {colour} in Figure \ref{fig08}. 
\begin{figure}[H]
\makebox[0.97\linewidth][r]{%
	\begin{subfigure}[b]{.68\textwidth}
	
			\includegraphics[width=0.95\textwidth]{Fig_08a.jpg}
		\caption{	\centering}
	\end{subfigure}%
	\begin{subfigure}[b]{.61\textwidth}
		
			\includegraphics[width=0.95\textwidth]{Fig_08b.jpg}
		\caption{\centering}
	\end{subfigure}%
} %MDPI: Please use commas to separate thousands for numbers with five or more digits (not four digits) in the picture, e.g., "10000" should be "10,000". <-- Answer: corrected (Figure 8a and 8b are replotted), including changed hyphen (-) into a minus sign ($-$).
       \caption{NDVI computed for multispectral {Landsat-9 OLI/TIRS} image.~({\bf a}) NDVI ratio-based vegetation index. %MDPI: We moved the subfigure explanations into the figure caption. Please confirm. <-- Answer: yes, we confirm.
 ({\bf b}) Data distribution within the computed NDVI.}
        \label{fig08}
\end{figure}


\subsubsection{Soil-Adjusted Vegetation Index (SAVI)}
The Soil-Adjusted Vegetation Index (SAVI) is a distance-based vegetation index which is based on the NDVI but includes an adjustment factor to reduce the noise from the NDVI by eliminating the effects from the canopy background and atmospheric conditions. Thus, it adjusts the NDVI for soil brightness where vegetation coverage is low in urban, densely built up areas and bare soil using the arithmetic expression in Equation (\ref{eq:2}), proposed in %MDPI: please confirm if * in equation should be correct to \times <--  Answer: yes, \times is correct.
~\cite{HUETE1988295}:

\begin{equation}\label{eq:2}
SAVI=\dfrac{(NIR-Red)}{(NIR+Red+L)}\times(1+L)
\end{equation}
where L is a soil brightness correction factor, which is defined as 0.5 to adjust to most land cover types. {Compared with the NDVI, }the SAVI{ }is a finer indicator, {which enables it to be }used {for the assessment of} { if }land coverage contains healthy green vegetation with correction for brightness, with the majority of values lying in the diapason from $-$0.1 to 0.1 (Figure \ref{fig09}). Thus, the main advantage of the SAVI is the emphasised colour of the vegetation canopy corrected for the soil brightness, because the SAVI excludes the atmospheric or meteorological effects and those from soil. Compared with the{ }NDVI and EVI, which were affected by the soil brightness, the SAVI{ demonstrated }better results in differentiating between{ }vegetation and soil{.}

\begin{figure}[H]
\makebox[0.97\linewidth][r]{%
	\begin{subfigure}[b]{.68\textwidth}
		
			\includegraphics[width=0.95\textwidth]{Fig_09a.jpg}
		\caption{\centering}
	\end{subfigure}%
	\begin{subfigure}[b]{.61\textwidth}
		
			\includegraphics[width=0.95\textwidth]{Fig_09b.jpg}
		\caption{\centering}
	\end{subfigure}%
}
       \caption{SAVI %MDPI: Please use commas to separate thousands for numbers with five or more digits (not four digits) in the picture, e.g., "10000" should be "10,000". <-- Answer: corrected (Figure 9a and 9b are replotted). 
 %MDPI: Please change the hyphen (-) into a minus sign ($-$, "U+2212"). e.g., "-1" should be "$-$1".  <-- Answer: corrected both in Figure 9a and 9b.
computed for multispectral image \emph{LC09\_L1TP\_197055\_20220111\_20220122\_02\_T1\_B}. ({\bf a}) SAVI ratio-based vegetation index. %MDPI: We moved the subfigure explanations into the figure caption. Please confirm. <-- Answer: yes, we confirm.
 ({\bf b}) Data distribution within the computed SAVI.}
        \label{fig09}
\end{figure}

Thus, the SAVI index was more suitable and better adjusted to agricultural monitoring. At the same time, because C\^ote d'Ivoire is a highly productive agricultural region, {being the} world's most important cocoa producer{, the }agricultural monitoring {of its landscapes} is of high importance{.~Therefore, }analysis of {the} vegetation indices is valuable for sustainable development and environmental planning in C\^ote d'Ivoire.

\subsubsection{Enhanced Vegetation Index 2 (EVI2)}

The EVI (Figure \ref{fig10}) is an optimised spectral imaging transformation of the NIR and red image bands aimed at enhancing the contribution of vegetation properties{. Computing this index is performed} using Equation (\ref{eq:3}), {originally proposed }by \cite{JIANG20083833}.


{EVI2} was developed {as an} updated and adjusted formula {of the EVI. In this form, it was developed} as a satellite vegetation index for the Moderate Resolution Imaging Spectroradiometers (MODIS) and adopted for the Landsat Thematic Mapper (TM) or Landsat Enhanced Thematic Mapper Plus (ETM+) series. The expression is defined as follows in Equation (\ref{eq:4}), which was optimised in \cite{JIANG20083833}: 

\begin{equation}\label{eq:4}
EVI2=2.4*\dfrac{(NIR-Red)}{(NIR+Red+1)}
\end{equation}

This allows a robust spatial and temporal correlation of the terrestrial chlorophyll and photosynthetic activity, which affects the structure of the canopy in various types of vegetation, including savannahs and forests. The ratios of the NIR and red spectral bands in the image enable decreasing the topographic effects from the Earth's relief. Therefore, the visualised EVI highlights and stresses smaller differences in the characteristics of spectral reflectance between vegetation and all other land cover types. The range of the values normally lies in the diapason between $-$0.2 and 0.2, as can be also seen in Figure \ref{fig10}:
 
\begin{equation}\label{eq:3}
EVI=\dfrac{(NIR-Red)}{(NIR+C1*Red-C2*Blue+L)}
\end{equation}

\begin{figure}[H]
\makebox[0.97\linewidth][r]{%
	\begin{subfigure}[b]{.68\textwidth}
		
			\includegraphics[width=0.95\textwidth]{Fig_10a.jpg}
		\caption{\centering}
	\end{subfigure}%
	\begin{subfigure}[b]{.61\textwidth}
	
			\includegraphics[width=0.95\textwidth]{Fig_10b.jpg}
		\caption{	\centering}
	\end{subfigure}%
}
       \caption{EVI2, %MDPI: Please use commas to separate thousands for numbers with five or more digits (not four digits) in the picture, e.g., "10000" should be "10,000". <-- Answer: corrected (Figure 10a and 10b are replotted). 
 %MDPI: Please change the hyphen (-) into a minus sign ($-$, "U+2212"). e.g., "-1" should be "$-$1". <-- Answer: corrected both in Figure 10a and 10b.
computed for multispectral image \emph{LC09\_L1TP\_197055\_20220111\_20220122\_02\_T1\_B}. ({\bf a}) EVI2 ratio-based vegetation index. %MDPI: We moved the subfigure explanations into the figure caption. Please confirm. <-- Answer: yes, we confirm.
 ({\bf b}) Data distribution within the computed EVI2.}
        \label{fig10}
\end{figure}
\subsubsection{Atmospherically Resistant Vegetation Index 2 (ARVI2)}

ARVI2 is adjusted to be resistant to atmospheric effects{. At} the same time, it is more sensitive to a wider spectrum of concentrations of chlorophyll in leaves. {Specifically, for} Côte d'Ivoire, the values of the ARVI ranged from $-$0.28 to $-$0.1 (Figure \ref{fig11}), although in general, the diapason of the ARVI values may vary from $-$1 to 1 by following Equation (\ref{eq:5}), {which was originally developed in} \cite{134076}: 

\begin{equation}\label{eq:5}
ARVI=\dfrac{(NIR-Red-y(Red-Blue))}{(NIR+Red-y(Red-Blue))}
\end{equation}

where y is the quotient derived from the components of atmospheric reflectance in the blue and red channel{.
The equation for the }adjusted ARVI (ARVI2), which is more sensitive to vegetation fraction, is {presented} in Equation (\ref{eq:6}), {which was originally developed in}~\cite{134076}: 

\begin{equation}\label{eq:6}
ARVI2=\dfrac{(NIR-Red)}{(NIR+Red)}*1.17 - 0.18
\end{equation}

This shows the regions of high atmospheric aerosol contents with higher moisture. Compared with the NDVI,{ }the ARVI adjusts and minimises the atmospheric scattering effects due to the blue light reflectance values, which also affects the reflectance in the red light \cite{134076}. 

Similar to{ }the NDVI, ARVI2 is also sensitive to the vegetation ratio in the landscape{ and }useful in climate and environmental studies{. }Specifically, {it can be applied }to evaluate the absorption of solar radiation, which is {highly} applicable in tropical African regions such as C\^ote d'Ivoire.

\begin{figure}[H]
\makebox[0.97\linewidth][r]{%
	\begin{subfigure}[b]{.68\textwidth}
	
			\includegraphics[width=0.95\textwidth]{Fig_11a.jpg}
		\caption{	\centering}
	\end{subfigure}%
	\begin{subfigure}[b]{.61\textwidth}
		
			\includegraphics[width=0.95\textwidth]{Fig_11b.jpg}
		\caption{\centering}
	\end{subfigure}%
}
       \caption{Atmospherically %MDPI: Please use commas to separate thousands for numbers with five or more digits (not four digits) in the picture, e.g., "10000" should be "10,000".  <-- Answer: corrected both in Figure 11a and 11b.
%MDPI: Please change the hyphen (-) into a minus sign ($-$, "U+2212"). e.g., "-1" should be "$-$1".  <-- Answer: corrected both in Figure 11a and 11b.
 Resistant Vegetation Index 2 (ARVI2) computed for multispectral image \emph{LC09\_L1TP\_197055\_20220111\_20220122\_02\_T1\_B}. ({\bf a}) ARVI2 ratio-based vegetation index. %MDPI: We moved the subfigure explanations into the figure caption. Please confirm. <-- Answer: yes, we confirm.
 ({\bf b}) Data distribution within the computed EVI2. }
        \label{fig11}
\end{figure}

%\newpage
\subsection{Terrain Analysis}

Cartographic visualisation is an essential application of satellite image analysis as a geospatial dataset. In this regard, terrain {mapping }was performed using{ }remote sensing data {and} the integrated approach of Python and R (Figure \ref{fig12}). 
\vspace{-6pt}
\begin{figure}[H]
%     \centering
\makebox[0.97\linewidth][r]{%
     \begin{subfigure}[r]{0.68\textwidth}
%         \centering
%        		\includegraphics[width=8.6cm]{Fig_12a}
		\includegraphics[width=0.95\textwidth]{Fig_12a}
         \caption{ \centering}
     \end{subfigure} \hfill \hspace{2mm}
     \begin{subfigure}[r]{0.56\textwidth}
 %        \centering
 %        \includegraphics[width=7cm]{Fig_12b}
         \includegraphics[width=0.95\textwidth]{Fig_12b}
         \caption{ \centering}
     \end{subfigure}
}
        \caption{(\textbf{a}): Landsat OLI/TIRS %MDPI: Please use commas to separate thousands for numbers with five or more digits (not four digits) in the picture, e.g., "10000" should be "10,000". <--  Answer: corrected (Figure 12a is replotted). 
 image with enlarged fragment of Kossou Lake, using grayscale visualisation of Band 6 (SWIR). Mapping: R. (\textbf{b}) Enlarged fragment of the cropped SRTM3 void-filled DEM overlayed on top of the hillshade. Cmaps: `turbo' and `gist\_yarg'. Mapping: Python.}
        \label{fig12}
\end{figure}

The data were linked with the digital elevation model (DEM) to demonstrate the effective approach of both{ } languages to modelling Earth observation data. {We selected} a DEM with a high spatial resolution {for terrain modelling because it }is one of the most crucial base layer factors for environmental modelling \cite{AVAND20221}{. }The advances in cartographic modelling using programming scripts in recent years changed the terrain visualisation from GUI-based GIS mapping to dynamic scripts with automatic representation of the output images. Processing the DEM from the SRTM using Python scripts based on the EarthPy library demonstrated the advantages and flexibility of the workflow, aimed at deriving topographic information for spatial analysis (Figure \ref{fig13}). The DEM developed from the SRTM of rugged terrain in Côte d'Ivoire{ had }a spatial resolution of 3 arc seconds, which corresponds to $\pm 90$ m.
\vspace{-6pt}
\begin{figure}[H]
\makebox[0.90\linewidth][c]{%
	\begin{subfigure}[b]{.6\textwidth}
		
			\includegraphics[width=0.9\textwidth]{Fig_13a}
		\caption{\centering}
	\end{subfigure}%
	\begin{subfigure}[b]{.6\textwidth}
		
			\includegraphics[width=0.9\textwidth]{Fig_13b}
		\caption{\centering}
	\end{subfigure}%
}\\
 \vfill 
\makebox[0.90\linewidth][c]{%
	\begin{subfigure}[b]{.6\textwidth}
		
			\includegraphics[width=0.9\textwidth]{Fig_13c}
		\caption{\centering}
	\end{subfigure}%
	\begin{subfigure}[b]{.6\textwidth}
	
			\includegraphics[width=0.9\textwidth]{Fig_13d}
		\caption{\centering}
	\end{subfigure}%
}
\caption{Terrain modelling for Kossou Lake area, Côte d'Ivoire. Mapping: Python. Source: authors. (\textbf{a}) DTM of Kossou Lake, %MDPI: We moved the subfigure explanations into the figure caption. Please confirm. <-- Answer: yes, we confirm.
 Côte d'Ivoire; 
 (\textbf{b}) Hillshade from DTM of Kossou Lake, Côte d'Ivoire. (\textbf{c})~Azimuth of Sun: 230\textdegree. Hillshade from DTM. (\textbf{d}) Sun angle altitude: 10\textdegree. Hillshade from DTM.} \label{fig13}
\end{figure}

The topographic variation was {modelled} with the Python {library} EarthPy{, }based on the processing of{ }fields (raster layers), spatial entities (elevated objects with 3D coordinates as attributes) and the network (topology of spatial links among the geographic objects). This approach {uses} relational data modelling with height as the primary identifier in a dataset. The DEM data are modelled by EarthPy using a raster data structure as a square grid matrix, with the topography defined by the sampled points and with XYZ dimensions indicating the continuous field of the spatial altitudes. 

{Spatial} modelling{ }continues topographic representation using discrete boundaries (Figure \ref{fig13}). {It }is based on the topographic attributes derived from the DEM and aims to{ }define {the geomorphic and hydrological} features {as} environmental gradients.~These derivatives include such attributes as the slope, aspect, upslope drainage and catchment areas, tangential, plan and profile curvature of the cross-sections, upslope and flow path length and hillshade. Determining the drainage direction and connectivity of the raster points {from} the DEM depends on{ }the surface attributes of the topographic hillslopes and stream channels which comprise the hydrological network system. In this regard,{ }morphometric modelling contributes to landscape analysis through visual representation of spatial variability and connectivity of the primary {terrain} attributes{ }distributed on the surface of the Earth. 

Numerous applications of topographic modelling in Earth and environmental studies motivated the development of Python programming approaches aimed at calculating the terrain. This is especially advantageous for the topographically diverse landscapes with interspersed forest and savannah areas which are{ }typical for{ }West Africa. The {principal} advantages of the EarthPy library consist of straightforward and robust topographic modelling, which presents the visualised series of the topographic attributes derived from a square grid DEM. These include the slope, aspect and topographic relief shading with{ }their attributes (i.e., {hillshade with varied }Sun azimuths and angle altitudes), as demonstrated in Figure \ref{fig13}.

The environmental effects from the terrain landforms are reflected in the soil characteristics, as its formation is regulated {by a variety of factors: }the geologic substrate,{ }direction of stream flow{ }and susceptibility of hillslopes to{ }soil erosion{. In turn, this} depends on{ }the slope steepness and relief curvature. Application of the terrain model includes analysis of the sediment transport, which is {controlled by the} terrain type and contributes to soil erosion. That aside, the morphometric aspects, differences {in} topographic heights, disposition of the slope and aspect and topographic horizon affect the distribution of solar radiation{. The latter} controls the soil vegetation system through{ }the biophysical processes of evapotranspiration and photosynthesis. Thus, the disposition of a hill slopes towards the sun angle, as the slope aspect contributes to the distribution of vegetation types through the effects of microclimate. For instance, more sensitive plants grow on the southern flanks of the hill slopes to benefit the solar energy. Thus, spatial continuity of the environmental processes is deeply linked to the regional topographic properties of the terrain.~Therefore, {modelling }the topographic attributes computed by the Python {library EarthPy, applied to a} high-resolution DEM, reveals{ }the variability of the terrain {as a} complex heterogeneous {structure} and contributes{ }to the environmental analysis of Côte d'Ivoire{, West Africa}.

%%%%%%%%%%%%%%%%%%%%%%%%%%%%%%%%%%%%%%%%%%
\section{Conclusions}

{In this paper, we proposed two scripting approaches to satellite image processing to compute several vegetation indices and to perform terrain modelling. The scripting techniques of the 'raster' and 'terra' packages of R for computing vegetation indices demonstrated great potential for using map algebra and calculation techniques with the Landsat OLI/TIRS for monitoring vegetation health and environmental mapping. Moreover, this technique is also useful for agricultural applications where deteriorated sites may require remediation or close monitoring. To model terrain variables, we proposed a Python-based approach where the topographic data are derived from the SRTM DEM using the EarthPy library.}

The problem of the terrain analysis in image processing can be referred {to }as spatial data modelling, with classification of cells according to the elevation values{. The derived} variables {are essential }for visualisation of the slope, aspect and hillshade using the parameters of the Sun's angle and azimuth.%Check meaning retained.
~Mapping the topographic data {using terrain analysis by Python }highlighted their fundamental importance in Earth modelling. {Thus, the }relief {of the Earth's surface }has {notable }effects on various components of the {ecosystems}, including{ }geological, climatic {and} hydrologic processes {as well as} the soil vegetation parameters. {Moreover, }topography regulates and controls climate processes {such as }precipitation, wind direction and strength. Therefore, it plays a crucial role in the environmental dynamics and ecosystem functioning, which includes the uneven distribution{ of the diverse landscapes with }contrasting types of vegetation in the hilly and coastal areas of Côte d'Ivoire. 

{The }integrated use of regional environmental resources presents {the} potential for sustainable development of Côte-d'Ivoire. This requires regular monitoring of the environmental conditions of the country using satellite images. In data science, processing satellite images using advanced programming methods is a challenge which consists of adaptation and optimisation of {the programming }algorithms towards spatial data modelling and analysis. In this study, we presented a solution to this problem with a straightforward yet effective approach which combines libraries of Python and R for processing {RS} data in a scripting framework. The use of both programming languages{ allowed us} to {use} the best packages for topographic and environmental modelling separately when handling the satellite images. {It also demonstrated the integration of the spatial} libraries of {the }two scripting languages{ }for{ }terrain mapping and{ }image processing {aimed at environmental data analysis}. 

The techniques of R {enabled} image processing{ }using information derived from pixels in the image scene. {Spatial data} were differentiated by the RGB triplet colours through plotting the colour composite of the Landsat 8 OLI/TIRS image using the R library 'raster', which {adapts} the general syntax of R for spatial data processing. The colours corresponded to the composites of the RGB signals scattered from the natural and built-up environment, which resulted in various DN values for the pixel cells in a raster matrix {of the image}. The combinations of the three channels were tested by the R library 'terra' as sensor bands, which {displayed} the image and {visualised} the appearance of various {terrain }objects{ }using the distinct colours. We compared the correlation between the selected bands of Landsat-8 OLI/TIRS C1 in a multispectral image for {Bands} 1 and 2 (blue and green) and 4 and 5 (red and NIR). We then computed four selected vegetation indices---NDVI, SAVI, EVI2 and ARVI2---using the principles of map algebra based on the DNs of the reflectance values of the pixels in the selected spectral bands (NIR and red). The major effect of spectral reflectance from healthy green vegetation was visible in the NIR region of the image spectrum, with a strong NIR response rising sharply for healthy vegetation. Conversely, the areas with poor or damaged vegetation were represented by a low NIR response, which notably decreased as shown by the vegetation indices and {with corresponding lower values for the NDVI}. 

{The algorithms of the EarthPy library of Python were used for terrain modelling based on the} high-resolution SRTM DEM. The upscaled mapping was shown at the two hierarchical levels: regional environmental analysis of the Yamoussoukro surroundings and local topographic modelling of the Kossou Lake. Processing the topographic data for{ }visualisation of the slope, aspect, hillshade and elevation {was based on} interpolation of the cell values {in a raster matrix grid} to spatial polygons{ that corresponded to the diverse terrain structures}. We {showed} these methods {in} a series of scripts with detailed explanations of their functionality and the effects{ of changed} parameters on the image. The main advantage of{ }Python for spatial modelling {consists of the machine-based automation of data processing. Specifically, it derives} information from the satellite images in an efficient way {using} a straightforward syntax, {which enables} handling complex {RS data}. {The }nonlinear environmental and topographic links among the near-surface processes {could be depicted from raster matrices of the satellite images} using programming algorithms. {This enabled us to highlight }the complex relationships between the objects forming the terrain shape and the diverse landscapes on Earth. 

A wide variety of Earth's features can be {mapped using} satellite images{: }climate change, urban sprawl, degradation {of vegetation}, deforestation, topographic variations{ }or land cover change.~At the same time, the analysis of the complexity of {a} particular{ }phenomenon requires advanced methods of data processing to reveal the links and correlations among the objects and to map spatial patterns. For instance, topographic variations{ }may be caused by geologic effects and anthropogenic activities, or the environmental dynamics of the Earth's landscapes may involve nonlinear long-term and short-term variations. The comprehensive analysis of such processes requires {advanced automatic} processing of large datasets, where a conventional GIS is not effective enough. {At the same time, due} to the high level of spatial detail, high-resolution satellite images require efficient algorithms for data handling, including image resampling, band composition, statistical processing and map algebra for calculation of the vegetation indices. This requires the advanced {methods for} image processing {made possible by programming languages}.

As we demonstrated{, }the effective way to handle this issue includes the {use} of high-level programming languages such as Python and R{. Scripting libraries are effective }for automatic image processing and retrieval of spatial information from {satellite} images. {The advanced }algorithms of {Python and R} can be successfully used for processing Earth observation data and geospatial modelling in a robust, automatic and effective way. {The application of }programming {approaches} for terrain analysis and environmental monitoring {is} successful in dealing with{ }multispectral {RS} data. The advantages of Python and R in image processing include{ }scripting algorithms which ensure automation of {RS} data processing.{ }Compared with{ }GISs{,} the functionality of Python and R{ }were proven{ }to have robust handling of the geospatial data{ due to the high }speed and effectiveness of code{ and repeatability of the scripts, which} outperformed{ }the existing{ state-of-the-art }approaches {of the traditional software}. 

In future works, we plan to {use} combinations of other libraries of Python and R for {higher-resolution images, such as} Sentinel-2A {or} SAR data for{ }environmental {mapping of} the tropical forest {regions} of Africa. Specifically, we intend to focus on the long-term analysis of the environmental dynamics using time series of satellite images. Spatial modelling of {the} large series of {RS data} involves{ }challenging problems with regard of their processing. For example, it is computationally challenging to deal with big data and requires complex methods of {data} handling. To this end, we plan to develop effective approaches to{ process }large spatiotemporal datasets using scripting languages.{ Our next work will continue the development of the advanced cartographic methods for processing satellite images with programming languages}.

\vspace{6pt}
\authorcontributions{Supervision, conceptualisation, methodology, software, resources, funding acquisition and project administration, O.D.; writing---original draft preparation, methodology, software, data curation, visualisation, formal analysis, validation, writing---review and editing and investigation, P.L. All authors have read and agreed to the published version of the manuscript.}

\funding{This project was supported by The Federal Public Planning Service Science Policy or Belgian Science Policy Office, Federal Science Policy - BELSPO (B2/202/P2/SEISMOSTORM).}

\institutionalreview{Not applicable.}

\informedconsent{Not applicable.}

\dataavailability{The open GitHub repository with Python and R scripts used for image processing in this study can be found at \url{https://github.com/paulinelemenkova/Cote_d_Ivoire_R_Python_Listings} (accessed on 2 November 2022).}

\acknowledgments{The authors thank the anonymous reviewers for their careful reading, critical suggestions and constructive comments that helped improve a previous version of this manuscript.}

\conflictsofinterest{The authors declare no conflict of interest.} 

\abbreviations{Abbreviations}{
The following abbreviations are used in this manuscript:\\

\noindent 
\begin{tabular}{@{}ll}
ARVI & ~~~~~~~~~~~~~~~~~~~~~~Atmospherically Resistant Vegetation Index \\
CPT & ~~~~~~~~~~~~~~~~~~~~~~Colour palette table \\
CRS & ~~~~~~~~~~~~~~~~~~~~~~Coordinate reference system \\
DCW & ~~~~~~~~~~~~~~~~~~~~~~Digital Chart of the World \\
DEM & ~~~~~~~~~~~~~~~~~~~~~~Digital elevation model \\
DN & ~~~~~~~~~~~~~~~~~~~~~~Digital number \\
DTM & ~~~~~~~~~~~~~~~~~~~~~~Digital terrain model \\
EVI &~~~~~~~~~~~~~~~~~~~~~~ Enhanced Vegetation Index \\
GEBCO & ~~~~~~~~~~~~~~~~~~~~~~General Bathymetric Chart of the Oceans \\
GDAL & ~~~~~~~~~~~~~~~~~~~~~~Geospatial Data Abstraction Library \\
GIS & ~~~~~~~~~~~~~~~~~~~~~~Geographic information system \\
GloVis & ~~~~~~~~~~~~~~~~~~~~~~Global Visualization Viewer \\
GMT & ~~~~~~~~~~~~~~~~~~~~~~Generic Mapping Tools \\
GPS & ~~~~~~~~~~~~~~~~~~~~~~Global positioning system \\
%GUI & Graphical User Interface \\
%Landsat 8-9 OLI/TIRS & Landsat 8-9 Operational Land Imager and Thermal Infrared Sensor \\
%Landsat TM & Landsat Thematic Mapper \\
%Landsat ETM+ & Landsat Enhanced Thematic Mapper Plus \\
%LiDAR & Light Detection and Ranging \\ 
%ML & Machine Learning \\
%MODIS & Moderate Resolution Imaging Spectroradiometers \\
%MNDWI & Modified Normalized Difference Water Index \\
%NASA & National Aeronautics and Space Administration \\
%NDVI & Normalized Difference Vegetation Index \\
%NDWI & Normalized Difference Water Index \\
%NGA & National Geospatial-Intelligence Agency \\
%NIR & Near Infrared \\
%RGB & Red Green Blue \\
%{RS} & {Remote Sensing} \\
%SAGA & System for Automated Geoscientific Analyses \\
%SAVI & Soil-Adjusted Vegetation Index \\
%SPOT & Satellite pour l'Observation de la Terre \\
%SRTM & Shuttle Radar Topography Mission \\
%SWIR & Short Wave InfraRed \\
%3D & Three-dimensional \\
%TIFF & Tagged Image File Format \\
%USGS & United States Geological Survey \\
%WRI & Water Ratio Index \\
\end{tabular}


\noindent 
\begin{tabular}{@{}ll}
%ARVI & Atmospherically Resistant Vegetation Index \\
%CPT & Colour Palette Table \\
%CRS & Coordinate Reference System \\
%DCW & Digital Chart of the World \\
%DEM & Digital Elevation Model \\
%DN & Digital Number \\
%DTM & Digital Terrain Model \\
%EVI & Enhanced Vegetation Index \\
%GEBCO & General Bathymetric Chart of the Oceans \\
%GDAL & Geospatial Data Abstraction Library \\
%GIS & Geographic Information System \\
%GloVis & Global Visualization Viewer \\
%GMT & Generic Mapping Tools \\
%GPS & Global Positioning System \\
GUI & Graphical user interface \\
Landsat 8-9 OLI/TIRS & Landsat 8-9 Operational Land Imager and Thermal Infrared Sensor \\
Landsat TM & Landsat Thematic Mapper \\
Landsat ETM+ & Landsat Enhanced Thematic Mapper Plus \\
LiDAR & Light detection and ranging \\ 
ML & Machine learning \\
MODIS & Moderate Resolution Imaging Spectroradiometers \\
MNDWI & Modified Normalized Difference Water Index \\
NASA & National Aeronautics and Space Administration \\
NDVI & Normalized Difference Vegetation Index \\
NDWI & Normalized Difference Water Index \\
NGA & National Geospatial-Intelligence Agency \\
NIR & Near infrared \\
RGB & Red, green and blue \\
{RS} & {Remote sensing} \\
SAGA & System for Automated Geoscientific Analyses \\
SAVI & Soil-Adjusted Vegetation Index \\
SPOT & Satellite pour l'Observation de la Terre \\
SRTM & Shuttle Radar Topography Mission \\
SWIR & Short-wave infrared \\
3D & Three-dimensional \\
TIFF & Tagged Image File Format \\
USGS & United States Geological Survey \\
WRI & Water Ratio Index \\
\end{tabular}
}

\begin{adjustwidth}{-\extralength}{0cm}
\reftitle{References}

% Please provide either the correct journal abbreviation (e.g. according to the “List of Title Word Abbreviations” http://www.issn.org/services/online-services/access-to-the-ltwa/) or the full name of the journal.
% Citations and References in Supplementary files are permitted provided that they also appear in the reference list here. 

%=====================================
% References, variant A: external bibliography
%=====================================
\begin{thebibliography}{999}

\bibitem[Coulibaly \em{et~al.}(2021)Coulibaly, Guan, Assoma, Fan, and
  Coulibaly]{COULIBALY2021108092}
Coulibaly, L.K.; Guan, Q.; Assoma, T.V.; Fan, X.; Coulibaly, N.
\newblock {Coupling linear spectral unmixing and RUSLE2 to model soil erosion
  in the Boubo coastal watershed, Côte d'Ivoire}.
\newblock {\em Ecol. Indic.} {\bf 2021}, {\em 130},~108092.
\newblock {https://doi.org/10.1016/j.ecolind.2021.108092}.

\bibitem[Koné \em{et~al.}(2016)Koné, Coulibaly, Kouadio, Neuba, and
  Malan]{7729337}
Koné, M.; Coulibaly, L.; Kouadio, Y.L.; Neuba, D.F.; Malan, D.F.
\newblock Multitemporal monitoring of the forest cover in Côte d'Ivoire from
  the 1960s to the 2000s, using Landsat satellite images.
\newblock In Proceedings of the 2016 IEEE International Geoscience and Remote
  Sensing Symposium (IGARSS), Beijing, China, 10--15 July 2016; pp. 1325--1328.
\newblock {https://doi.org/10.1109/IGARSS.2016.7729337}. %MDPI: newly added information, please confirm. <-- Answer: yes, we confirm.

\bibitem[Lemenkova(2020)]{Lemenkova2020d}
Lemenkova, P.
\newblock {Sentinel-2 for High Resolution Mapping of Slope-Based Vegetation
  Indices Using Machine Learning by SAGA GIS}.
\newblock {\em Transylv. Rev. Syst. Ecol. Res.}
  {\bf 2020}, {\em 22},~17--34.
\newblock {https://doi.org/10.2478/trser-2020-0015}.

\bibitem[Wang \em{et~al.}(2014)Wang, Wang, and Li]{6946534}
Wang, J.; Wang, H.; Li, X.
\newblock Down scaling vegetation fraction by fusing multi-temporal MODIS and
  Landsat data.
\newblock In Proceedings of the 2014 IEEE Geoscience and Remote Sensing
  Symposium, Quebec City, QC, USA, 13--18 July 2014; pp. 757--760.
\newblock {https://doi.org/10.1109/IGARSS.2014.6946534}. %MDPI: newly added information, please confirm. <-- Answer: yes, we confirm.

\bibitem[Lemenkova(2020)]{Lemenkova202019}
Lemenkova, P.
\newblock {Hyperspectral Vegetation Indices Calculated by Qgis Using Landsat Tm
  Image: a Case Study of Northern Iceland}.
\newblock {\em Adv. Res. Life Sci.} {\bf 2020}, {\em 4},~70--78.
\newblock {https://doi.org/10.2478/arls-2020-0021}.

\bibitem[Wang \em{et~al.}(2012)Wang, Hajnsek, and Kinzelbach]{6352694}
Wang, H.; Hajnsek, I.; Kinzelbach, W.
\newblock Calibrated Landsat TM LAI retrieval for monitoring vegetation cover
  change after ecological releases to the lower Tarim river.
\newblock In Proceedings of the 2012 IEEE International Geoscience and Remote
  Sensing Symposium, Munich, Germany, 22--27 July 2012; pp. 6279--6282.
\newblock {https://doi.org/10.1109/IGARSS.2012.6352694}. %MDPI: newly added information, please confirm. <-- Answer: yes, we confirm.

\bibitem[N’go(2015)]{2015hydrologie}
N’go, Y.
\newblock Hydrologie et dynamique de l’{\'e}tat de surface des terres dans le
  sud-ouest de la C{\^o}te d’Ivoire: Impacts et moteurs de d{\'e}gradation.
\newblock {\em Doctorat d’{\'e}tat {\`e}s Sciences Naturelles de
  l’Universit{\'e} Nangui Abrogoua en Sciences et Gestion de
  l’Environnemen, C{\^o}te d’Ivoire}, \hl{2015}. %MDPI: Please provide more information about the article type, such as book (please provide the name and location of the publisher); online resource (please provide the URL of the website and the date it was accessed (Date Month Year)); or journal article (please provide the name of the journal, the year and volume in which it was published, and the page number). Please refer to https://www.mdpi.com/authors/references for full reference formatting guides.

\bibitem[Lees \em{et~al.}(1985)Lees, Lettis, and Bernstein]{1457509}
Lees, R.D.; Lettis, W.R.; Bernstein, R.
\newblock Evaluation of Landsat thematic mapper imagery for geologic
  applications.
\newblock {\em Proc.  IEEE} {\bf 1985}, {\em 73},~1108--1117.
\newblock {https://doi.org/10.1109/PROC.1985.13241}.

\bibitem[Ngom \em{et~al.}(2022)Ngom, Mbaye, Baratoux, Baratoux, Ahoussi,
  Kouame, Faye, and Sow]{NGOM2022102873}
Ngom, N.M.; Mbaye, M.; Baratoux, D.; Baratoux, L.; Ahoussi, K.E.; Kouame, J.K.;
  Faye, G.; Sow, E.H.
\newblock {Recent expansion of artisanal gold mining along the Bandama River
  (Côte d’Ivoire)}.
\newblock {\em Int. J. Appl. Earth Obs. Geoinf.} {\bf 2022}, {\em 112},~102873.
\newblock {https://doi.org/10.1016/j.jag.2022.102873}.

\bibitem[Jessell \em{et~al.}(2015)Jessell, Santoul, Baratoux, Youbi, Ernst,
  Metelka, Miller, and Perrouty]{JESSELL2015440}
Jessell, M.; Santoul, J.; Baratoux, L.; Youbi, N.; Ernst, R.E.; Metelka, V.;
  Miller, J.; Perrouty, S.
\newblock {An updated map of West African mafic dykes}.
\newblock {\em J. Afr. Earth Sci.} {\bf 2015}, {\em
  112},~440--450.
\newblock %MDPI: please confirm if this sentence could be deleted. <-- Answer: yes, we confirm (deleted).
{https://doi.org/10.1016/j.jafrearsci.2015.01.007}.

\bibitem[Siagné \em{et~al.}(2022)Siagné, Aïfa, Kouamelan, Houssou, Digbeu,
  Kakou, and Couderc]{SIAGNE2022104680}
Siagné, Z.H.; Aïfa, T.; Kouamelan, A.N.; Houssou, N.N.; Digbeu, W.; Kakou,
  B.K.F.; Couderc, P.
\newblock New lithostructural map of the Doropo region, northeast Côte
  d'Ivoire: Insight from structural and aeromagnetic data.
\newblock {\em J. Afr. Earth Sci.} {\bf 2022}, {\em
  196},~104680.
\newblock {https://doi.org/10.1016/j.jafrearsci.2022.104680}.

\bibitem[Sharma \em{et~al.}(2021)Sharma, Garg, Badenko, Fedotov, Min, and
  Yao]{SHARMA2021317}
Sharma, M.; Garg, R.D.; Badenko, V.; Fedotov, A.; Min, L.; Yao, A.
\newblock Potential of airborne LiDAR data for terrain parameters extraction.
\newblock {\em Quat. Int.} {\bf 2021}, {\em 575-576},~317--327.
\newblock %MDPI: please confirm if this sentence should be deleted <-- Answer: yes, we confirm (delted).
{https://doi.org/10.1016/j.quaint.2020.07.039}.

\bibitem[Kulkarni \em{et~al.}(2019)Kulkarni, Chandrashekaraiah, and S]{8978301}
Kulkarni, S.; Chandrashekaraiah, M.S.R.
\newblock {3D Annotation Tool Using LiDAR}.
\newblock In Proceedings of the 2019 Global Conference for Advancement in
  Technology (GCAT), Bangaluru, India, 18--20 October 2019; pp. 1--4.
\newblock {https://doi.org/10.1109/GCAT47503.2019.8978301}. %MDPI: newly added information, please confirm. <-- Answer: yes, we confirm.

\bibitem[Moffatt \em{et~al.}(2020)Moffatt, Platt, Mondragon, Kwok, Uryeu, and
  Bhandari]{9213897}
Moffatt, A.; Platt, E.; Mondragon, B.; Kwok, A.; Uryeu, D.; Bhandari, S.
\newblock {Obstacle Detection and Avoidance System for Small UAVs using a
  LiDAR}.
\newblock In Proceedings of the 2020 International Conference on Unmanned
  Aircraft Systems (ICUAS), Athens, Greece, 1--4 September 2020; pp. 633--640.
\newblock {https://doi.org/10.1109/ICUAS48674.2020.9213897}. %MDPI: newly added information, please confirm. <-- Answer: yes, we confirm.

\bibitem[Yeh \em{et~al.}(2017)Yeh, Lin, Chan, Chang, and Hsieh]{YEH2017236}
Yeh, C.H.; Lin, M.L.; Chan, Y.C.; Chang, K.J.; Hsieh, Y.C.
\newblock {Dip-slope mapping of sedimentary terrain using polygon auto-tracing
  and airborne LiDAR topographic data}.
\newblock {\em Eng. Geol.} {\bf 2017}, {\em 222},~236--249.
\newblock {https://doi.org/10.1016/j.enggeo.2017.04.009}.

\bibitem[Montreuil \em{et~al.}(2021)Montreuil, Chen, Moelans, Dierckx,
  Houthuys, Klein, and Bogaert]{9553919}
Montreuil, A.L.; Chen, M.; Moelans, R.; Dierckx, W.; Houthuys, R.; Klein, A.P.;
  Bogaert, P.
\newblock Monitoring Intertidal Bars and 3D Coastal Mapping Using an Automatic
  Algorithm on a Lidar Dataset.
\newblock In Proceedings of the 2021 IEEE International Geoscience and Remote
  Sensing Symposium IGARSS, Brussels, Belgium, 11--16 July 2021; pp. 8604--8607.
\newblock {https://doi.org/10.1109/IGARSS47720.2021.9553919}. %MDPI: newly added information, please confirm. <-- Answer: yes, we confirm.

\bibitem[Li \em{et~al.}(2019)Li, Xu, Macrander, Atkinson, Thomas, and
  Lopez]{LI201955}
Li, J.; Xu, Y.; Macrander, H.; Atkinson, L.; Thomas, T.; Lopez, M.A.
\newblock {GPU-based lightweight parallel processing toolset for LiDAR data for
  terrain analysis}.
\newblock {\em Environ. Model. Softw.} {\bf 2019}, {\em
  117},~55--68.
\newblock {https://doi.org/10.1016/j.envsoft.2019.03.014}.

\bibitem[Gupta and Sharma(2022)]{GUPTA2022100817}
Gupta, R.; Sharma, L.K.
\newblock {Mixed tropical forests canopy height mapping from spaceborne LiDAR
  GEDI and multisensor imagery using machine learning models}.
\newblock {\em Remote Sens. Appl. Soc. Environ.} {\bf
  2022}, {\em 27},~100817.
\newblock {https://doi.org/10.1016/j.rsase.2022.100817}.

\bibitem[Yuan \em{et~al.}(2022)Yuan, Choi, and Bolkas]{YUAN2022106966}
Yuan, W.; Choi, D.; Bolkas, D.
\newblock {GNSS-IMU-assisted colored ICP for UAV-LiDAR point cloud registration
  of peach trees}.
\newblock {\em Comput. Electron. Agric.} {\bf 2022}, {\em
  197},~106966.
\newblock {https://doi.org/10.1016/j.compag.2022.106966}.

\bibitem[Singhai and Rawat(2007)]{4426419}
Singhai, J.; Rawat, P.
\newblock Image enhancement method for underwater, ground and satellite images
  using brightness preserving histogram equalization with maximum entropy.
\newblock In Proceedings of the International Conference on Computational
  Intelligence and Multimedia Applications (ICCIMA 2007), Sivakasi,  India, 13--15 December 2007; Volume~3, pp. 507--512.
\newblock {https://doi.org/10.1109/ICCIMA.2007.359}. %MDPI: newly added information, please confirm. <-- Answer: yes, we confirm.

\bibitem[Yasuda \em{et~al.}(1993)Yasuda, Yamashita, Ruan, and Lu]{322052}
Yasuda, A.; Yamashita, K.; Ruan, Z.; Lu, Y.
\newblock Evaluation of brightness resolution of WEFAX images transmitted
  through GMS.
\newblock In Proceedings of the IGARSS '93---IEEE International
  Geoscience and Remote Sensing Symposium, Tokyo, Japan, 18--21 August 1993; Volume 4, pp. 1969--1971 .
\newblock {https://doi.org/10.1109/IGARSS.1993.322052}. %MDPI: newly added information, please confirm. <-- Answer: yes, we confirm.

\bibitem[Prudyus \em{et~al.}(2010)Prudyus, Lazko, and Semenov]{5446016}
Prudyus, I.; Lazko, L.; Semenov, S.
\newblock Satellite images quality improvement for multilevel data processing.
\newblock In Proceedings of the 2010 International Conference on Modern
  Problems of Radio Engineering, Telecommunications and Computer Science
  (TCSET), Lviv, Ukraine, 23--27 February 2010; pp. 56--58. %MDPI: newly added information, please confirm. <-- Answer: yes, we confirm.

\bibitem[Kita(2008)]{4761020}
Kita, Y.
\newblock A study of change detection from satellite images using joint
  intensity histogram.
\newblock In Proceedings of the 2008 19th International Conference on Pattern
  Recognition, Piscataway, NJ, USA, 8--11 December 2008; pp. 1--4.
\newblock {https://doi.org/10.1109/ICPR.2008.4761020}. %MDPI: newly added information, please confirm. <-- Answer: yes, we confirm.

\bibitem[Suresh and Hovenbitzer(2018)]{8519176}
Suresh, G.; Hovenbitzer, M.
\newblock {Texture and Intensity Based Land Cover Classification in Germany
  from Multi-Orbit \& Multi-Temporal Sentinel-1 Images}.
\newblock In Proceedings of the IGARSS 2018---2018 IEEE International
  Geoscience and Remote Sensing Symposium, Valencia, Spain, 22--27 July 2018; pp. 826--829.
\newblock {https://doi.org/10.1109/IGARSS.2018.8519176}. %MDPI: newly added information, please confirm. <-- Answer: yes, we confirm.

\bibitem[Yang and Hung(2002)]{995559}
Yang, S.; Hung, C.C.
\newblock Texture classification in remotely sensed images.
\newblock In Proceedings of the IEEE SoutheastCon 2002 (Cat.
  No.02CH37283), Valencia, Spain, 5--7 April 2002; pp. 62--66.
\newblock {https://doi.org/10.1109/SECON.2002.995559}. %MDPI: newly added information, please confirm. <-- Answer: yes, we confirm.

\bibitem[Lakshmanan \em{et~al.}(2000)Lakshmanan, DeBrunner, and Rabin]{899813}
Lakshmanan, V.; DeBrunner, V.; Rabin, R.
\newblock Texture-based segmentation of satellite weather imagery.
\newblock In Proceedings of the 2000 International Conference on
  Image Processing (Cat. No.00CH37101), Vancouver, BC, Canada, 10--13 September 2000%MDPI: Please add the location of the conference. <-- Answer: added (Vancouver, BC, Canada).
; Volume~2, pp. 732--735.
\newblock {https://doi.org/10.1109/ICIP.2000.899813}.

\bibitem[Jun and Tingting(2019)]{8868688}
Jun, X.; Tingting, S.
\newblock Study on Super-Resolution of Images Obtained by Micro Satellite with
  CMOS Sensor.
\newblock In Proceedings of the 2019 IEEE 4th International Conference on
  Signal and Image Processing (ICSIP), Wuxi, China, 19--21 July 2019; pp. 907--910.
\newblock {https://doi.org/10.1109/SIPROCESS.2019.8868688}. %MDPI: newly added information, please confirm. <-- Answer: yes, we confirm.

\bibitem[Perin \em{et~al.}(2021)Perin, Tulbure, Gaines, Reba, and
  Yaeger]{PERIN2021106694}
Perin, V.; Tulbure, M.G.; Gaines, M.D.; Reba, M.L.; Yaeger, M.A.
\newblock On-farm reservoir monitoring using Landsat inundation datasets.
\newblock {\em Agric. Water Manag.} {\bf 2021}, {\em 246},~106694.
\newblock {https://doi.org/10.1016/j.agwat.2020.106694}.

\bibitem[Pahlevan \em{et~al.}(2022)Pahlevan, Smith, Alikas, Anstee, Barbosa,
  Binding, Bresciani, Cremella, Giardino, Gurlin, Fernandez, Jamet, Kangro,
  Lehmann, Loisel, Matsushita, Hà, Olmanson, Potvin, Simis, VanderWoude,
  Vantrepotte, and Ruiz-Verdù]{PAHLEVAN2022112860}
Pahlevan, N.; Smith, B.; Alikas, K.; Anstee, J.; Barbosa, C.; Binding, C.;
  Bresciani, M.; Cremella, B.; Giardino, C.; Gurlin, D.;  et~al.
\newblock Simultaneous retrieval of selected optical water quality indicators
  from Landsat-8, Sentinel-2, and Sentinel-3.
\newblock {\em Remote Sens. Environ.} {\bf 2022}, {\em 270},~112860.
\newblock {https://doi.org/10.1016/j.rse.2021.112860}.

\bibitem[Lemenkova(2021)]{Lemenkova202141}
Lemenkova, P.
\newblock {Robust Vegetation Detection Using RGB Colour Composites and Isoclust
  Classification of the Landsat TM Image}.
\newblock {\em Geomat. Landmanag. Landsc.} {\bf 2021}, {\em
  4},~147--167.
\newblock {https://doi.org/10.15576/GLL/2021.4.147}.

\bibitem[Abu \em{et~al.}(2021)Abu, Szantoi, Brink, Robuchon, and
  Thiel]{ABU2021107863}
Abu, I.O.; Szantoi, Z.; Brink, A.; Robuchon, M.; Thiel, M.
\newblock Detecting cocoa plantations in Côte d’Ivoire and Ghana and their
  implications on protected areas.
\newblock {\em Ecol. Indic.} {\bf 2021}, {\em 129},~107863.
\newblock {https://doi.org/10.1016/j.ecolind.2021.107863}.

\bibitem[Attoumane \em{et~al.}(2022)Attoumane, Stéphanie, Kacou, André,
  Karamoko, Seguis, and Zahiri]{ATTOUMANE2022103285}
Attoumane, A.; Stéphanie, D.S.; Kacou, M.; André, A.D.; Karamoko, A.W.;
  Seguis, L.; Zahiri, E.P.
\newblock Individual perceptions on rainfall variations versus precipitation
  trends from satellite data: An interdisciplinary approach in two
  socio-economically and topographically contrasted districts in Abidjan, Côte
  d'Ivoire.
\newblock {\em Int. J. Disaster Risk Reduct.} {\bf 2022},
  {\em 81},~103285.
\newblock {https://doi.org/10.1016/j.ijdrr.2022.103285}.

\bibitem[Masolele \em{et~al.}(2021)Masolele, {De Sy}, Herold, Marcos,
  Verbesselt, Gieseke, Mullissa, and Martius]{MASOLELE2021112600}
Masolele, R.N.; {De Sy}, V.; Herold, M.; Marcos, D.; Verbesselt, J.; Gieseke,
  F.; Mullissa, A.G.; Martius, C.
\newblock Spatial and temporal deep learning methods for deriving land-use
  following deforestation: A pan-tropical case study using Landsat time series.
\newblock {\em Remote Sens. Environ.} {\bf 2021}, {\em 264},~112600.
\newblock {https://doi.org/10.1016/j.rse.2021.112600}.

\bibitem[Luo and Xu(2019)]{9066075}
Luo, X.; Xu, S.
\newblock {Forest Mapping from Hyperspectral Image Using Deep Belief Network}.
\newblock In Proceedings of the 2019 15th International Conference on Mobile
  Ad-Hoc and Sensor Networks (MSN),  Shenzhen, China, 11--13 December 2019; pp. 395--398.
\newblock {https://doi.org/10.1109/MSN48538.2019.00081}.

\bibitem[Martins \em{et~al.}(2022)Martins, Roy, Huang, Boschetti, Zhang, and
  Yan]{MARTINS2022113203}
Martins, V.S.; Roy, D.P.; Huang, H.; Boschetti, L.; Zhang, H.K.; Yan, L.
\newblock Deep learning high resolution burned area mapping by transfer
  learning from Landsat-8 to PlanetScope.
\newblock {\em Remote Sens. Environ.} {\bf 2022}, {\em 280},~113203.
\newblock {https://doi.org/10.1016/j.rse.2022.113203}.

\bibitem[Balzter \em{et~al.}(2015)Balzter, Cole, Thiel, and
  Schmullius]{rs71114876}
Balzter, H.; Cole, B.; Thiel, C.; Schmullius, C.
\newblock Mapping CORINE Land Cover from Sentinel-1A SAR and SRTM Digital
  Elevation Model Data using Random Forests.
\newblock {\em Remote Sens.} {\bf 2015}, {\em 7},~14876--14898.
\newblock {https://doi.org/10.3390/rs71114876}.

\bibitem[Dobesova and Dobes(2012)]{6344536}
Dobesova, Z.; Dobes, P.
\newblock {Comparison of visual languages in Geographic Information Systems}.
\newblock In Proceedings of the 2012 IEEE Symposium on Visual Languages and
  Human-Centric Computing (VL/HCC), Innsbruck, Austria, 30 September--4 October 2012; pp. 245--246.
\newblock {https://doi.org/10.1109/VLHCC.2012.6344536}.

\bibitem[Ellefsen \em{et~al.}(2019)Ellefsen, Lock, Settle, Karsten, and
  Parker]{ELLEFSEN20191171}
Ellefsen, K.L.; Lock, J.T.; Settle, B.; Karsten, C.A.; Parker, I.
\newblock {Applications of FLIKA, a Python-based image processing and analysis
  platform, for studying local events of cellular calcium signaling}.
\newblock {\em Biochim. Biophys. Acta (Bba) Mol. Cell Res.}
  {\bf 2019}, {\em 1866},~1171--1179.
\newblock %MDPI: Please check if this could be removed. <-- Answer: yes, we confirm (removed).
 {https://doi.org/10.1016/j.bbamcr.2018.11.012}.

\bibitem[Lemenkova(2022)]{Lemenkova202212}
Lemenkova, P.
\newblock {Handling Dataset with Geophysical and Geological Variables on the
  Bolivian Andes by the GMT Scripts}.
\newblock {\em Data} {\bf 2022}, {\em 7},~74.
\newblock {https://doi.org/10.3390/data7060074}.

\bibitem[Zhang \em{et~al.}(2014)Zhang, Yue, and Guo]{6910592}
Zhang, M.; Yue, P.; Guo, X.
\newblock GIScript: Towards an interoperable geospatial scripting language for
  GIS programming.
\newblock In Proceedings of the 2014 The Third International Conference on
  Agro-Geoinformatics, Beijing, China, 11--14 August 2014; pp. 1--5.
\newblock {https://doi.org/10.1109/Agro-Geoinformatics.2014.6910592}.  %MDPI: newly added information, please confirm. <-- Answer: yes, we confirm.

\bibitem[Lemenkova(2022)]{Lemenkova202215}
Lemenkova, P.
\newblock {Cartographic scripts for seismic and geophysical mapping of
  Ecuador}.
\newblock {\em Geografie} {\bf 2022}, {\em 127},~1--24.
\newblock {https://doi.org/10.37040/geografie.2022.006}.

\bibitem[Shi(2007)]{4160259}
Shi, X.
\newblock {Python for Internet GIS Applications}.
\newblock {\em Comput. Sci. Eng.} {\bf 2007}, {\em
  9},~56--59.
\newblock {https://doi.org/10.1109/MCSE.2007.57}.

\bibitem[De~Sarkar \em{et~al.}(2012)De~Sarkar, Biyahut, Kritika, and
  Singh]{6407912}
De~Sarkar, A.; Biyahut, N.; Kritika, S.; Singh, N.
\newblock {An environment monitoring interface using GRASS GIS and Python}.
\newblock In Proceedings of the 2012 Third International Conference on Emerging
  Applications of Information Technology, Kolkata, India, 16--18 April 2012; pp. 235--238.
\newblock {https://doi.org/10.1109/EAIT.2012.6407912}.  %MDPI: newly added information, please confirm. <-- Answer: yes, we confirm.

\bibitem[Pan \em{et~al.}(2015)Pan, Yang, and Xie]{7378577}
\textls[-25]{Pan, Z.; Yang, X.; Xie, Z.
\newblock A middleware: Python plugin transform on different GIS platforms.
\newblock In Proceedings of the 2015 23rd International Conference on
  Geoinformatics, Wuhan, China, 19--21 June 2015; pp. 1--6.
\newblock {https://doi.org/10.1109/GEOINFORMATICS.2015.7378577}.}  %MDPI: newly added information, please confirm. <-- Answer: yes, we confirm.

\bibitem[Vos \em{et~al.}(2019)Vos, Splinter, Harley, Simmons, and
  Turner]{VOS2019104528}
Vos, K.; Splinter, K.D.; Harley, M.D.; Simmons, J.A.; Turner, I.L.
\newblock CoastSat: A Google Earth Engine-enabled Python toolkit to extract
  shorelines from publicly available satellite imagery.
\newblock {\em Environ. Model. Softw.} {\bf 2019}, {\em
  122},~104528.
\newblock {https://doi.org/10.1016/j.envsoft.2019.104528}.

\bibitem[Van~Rossum and Drake~Jr(1995)]{van1995python}
Van~Rossum, G.; Drake, F.L., Jr.
\newblock {\em Python Reference Manual}; Centrum voor Wiskunde en Informatica:
  Amsterdam, The Netherlands, 1995.

\bibitem[Silva \em{et~al.}(2003)Silva, Lotufo, Machado, and Saude]{1247428}
Silva, A.; Lotufo, R.; Machado, R.; Saude, A.
\newblock Toolbox of image processing using the Python language.
\newblock In Proceedings of the 2003 International Conference on
  Image Processing (Cat. No.03CH37429), Seville, Spain, 14--17 September 2003; Volume~3, p. 1049.
\newblock {https://doi.org/10.1109/ICIP.2003.1247428}. %MDPI: newly added information, please confirm. <-- Answer: yes, we confirm.

\bibitem[Rey and Anselin(2010{\natexlab{a}})]{Rey2010PySALAP}
Rey, S.J.; Anselin, L.
\newblock PySAL: A Python Library of Spatial Analytical Methods.
\newblock {\em  Rev. Reg. Stud.} {\bf 2010}, {\em 37},~5--27.
\newblock {https://doi.org/10.52324/001c.8285}.

\bibitem[Rey and Anselin(2010{\natexlab{b}})]{Rey2010}
Rey, S.J.; Anselin, L., PySAL: A Python Library of Spatial Analytical Methods.
\newblock In {\em Handbook of Applied Spatial Analysis: Software Tools, Methods
  and Applications}; Springer: Berlin/Heidelberg, Germany, 2010; pp. 175--193.
\newblock {https://doi.org/10.1007/978-3-642-03647-7\_11}.

\bibitem[{de Deus Filho} \em{et~al.}(2022){de Deus Filho}, {da Silva Nunes},
  and Xavier]{DEDEUSFILHO2022101131}
{de Deus Filho}, J.C.A.; {da Silva Nunes}, L.C.; Xavier, J.M.C.
\newblock {iCorrVision-2D: An integrated python-based open-source Digital Image
  Correlation software for in-plane measurements (Part 1)}.
\newblock {\em SoftwareX} {\bf 2022}, {\em 19},~101131.
\newblock {https://doi.org/10.1016/j.softx.2022.101131}.

\bibitem[Farrens \em{et~al.}(2020)Farrens, Grigis, {El Gueddari}, Ramzi, G.R.,
  Starck, Sarthou, Cherkaoui, Ciuciu, and Starck]{FARRENS2020100402}
Farrens, S.; Grigis, A.; {El Gueddari}, L.; Ramzi, Z.; G.R., C.; Starck, S.;
  Sarthou, B.; Cherkaoui, H.; Ciuciu, P.; Starck, J.L.
\newblock {PySAP: Python Sparse Data Analysis Package for multidisciplinary
  image processing}.
\newblock {\em Astron. Comput.} {\bf 2020}, {\em 32},~100402.
\newblock {https://doi.org/10.1016/j.ascom.2020.100402}.

\bibitem[Gao \em{et~al.}(2022)Gao, Bilskie, and Hagen]{GAO2022105503}
Gao, S.; Bilskie, M.V.; Hagen, S.C.
\newblock {PyVF: A python program for extracting vertical features from
  LiDAR-DEMs}.
\newblock {\em Environ. Model. Softw.} {\bf 2022}, {\em
  157},~105503.
\newblock {https://doi.org/10.1016/j.envsoft.2022.105503}.

\bibitem[Dănilă \em{et~al.}(2012)Dănilă, Cazacu, and Gurlui]{6528254}
Dănilă, M.N.; Cazacu, M.M.; Gurlui, S.
\newblock Python utility: Laser-atmosphere interaction extended to network data
  management.
\newblock In Proceedings of the 2012 5th Romania Tier 2 Federation Grid, Cloud
  \& High Performance Computing Science (RQLCG), Cluj-Napoca, Romania, 25 October 2012; pp. 90--92;
\newblock {INSPEC Accession No: 13579885}. %MDPI: newly added information, please confirm. <-- Answer: yes, we confirm.

\bibitem[Roberts \em{et~al.}(2022)Roberts, Mwangi, Mukabi, Njui, Nzioka,
  Ndambiri, Bispo, Espirito-Santo, Gou, Johnson, Louis, Pacheco-Pascagaza,
  Rodriguez-Veiga, Tansey, Upton, Robb, and Balzter]{ROBERTS2022105192}
Roberts, J.F.; Mwangi, R.; Mukabi, F.; Njui, J.; Nzioka, K.; Ndambiri, J.K.;
  Bispo, P.C.; Espirito-Santo, F.D.B.; Gou, Y.; Johnson, S.C.M.;  et~al.
\newblock {Pyeo: A Python package for near-real-time forest cover change
  detection from Earth observation using machine learning}.
\newblock {\em Comput. Geosci.} {\bf 2022}, {\em 167},~105192.
\newblock {https://doi.org/10.1016/j.cageo.2022.105192}.

\bibitem[Chen \em{et~al.}(2022)Chen, Judge, and Hulse]{CHEN2022105362}
Chen, C.; Judge, J.; Hulse, D.
\newblock {PyLUSAT: An open-source Python toolkit for GIS-based land use
  suitability analysis}.
\newblock {\em Environ. Model. Softw.} {\bf 2022}, {\em
  151},~105362.
\newblock {https://doi.org/10.1016/j.envsoft.2022.105362}.

\bibitem[Wasser \em{et~al.}(2019)Wasser, Joseph, McGlinchy, Palomino, Korinek,
  Holdgraf, and Head]{wasser2019EarthPyAP}
Wasser, L.; Joseph, M.B.; McGlinchy, J.; Palomino, J.; Korinek, N.; Holdgraf,
  C.; Head, T.D.
\newblock EarthPy: A Python package that makes it easier to explore and plot
  raster and vector data using open source Python tools.
\newblock {\em J. Open Source Softw.} {\bf 2019}, {\em 4},~1886.
\newblock {https://doi.org/10.21105/joss.01886}.

\bibitem[Stančin and Jović(2019)]{8757088}
Stančin, I.; Jović, A.
\newblock An overview and comparison of free Python libraries for data mining
  and big data analysis.
\newblock In Proceedings of the 2019 42nd International Convention on
  Information and Communication Technology, Electronics and Microelectronics
  (MIPRO), Opatija, Croatia, 20--24 May 2019; pp. 977--982. %MDPI: newly added information, please confirm. <-- Answer: yes, we confirm.
\newblock {https://doi.org/10.23919/MIPRO.2019.8757088}.

\bibitem[Debeir and Decaestecker(2019)]{DebeirDecaestecker}
Debeir, O.; Decaestecker, C.
\newblock Data augmentation for training deep regression for in vitro cell
  detection.
\newblock In Proceedings of the 2019 Fifth International Conference on Advances
  in Biomedical Engineering (ICABME), Tripoli, Lebanon, 17--19 October 2019; pp. 1--3.
\newblock {https://doi.org/10.1109/ICABME47164.2019.8940275}. %MDPI: newly added information, please confirm. <-- Answer: yes, we confirm.

\bibitem[Debeir \em{et~al.}(2008)Debeir, Adanja, Warzee, Van~Ham, and
  Decaestecker]{Debeir2008}
Debeir, O.; Adanja, I.; Warzee, N.; Van~Ham, P.; Decaestecker, C.
\newblock Phase contrast image segmentation by weak watershed transform
  assembly.
\newblock In Proceedings of the 2008 5th IEEE International Symposium on
  Biomedical Imaging: From Nano to Macro, Paris, France, 14--17 May 2008; pp. 724--727.
\newblock {https://doi.org/10.1109/ISBI.2008.4541098}. %MDPI: newly added information, please confirm. <-- Answer: yes, we confirm.

\bibitem[Vangara \em{et~al.}(2021)Vangara, Vuddanti, and Kakani]{9719489}
Vangara, V.K.M.; Vuddanti, S.; Kakani, B.
\newblock An Accurate and Fast Computational Python Based Module for Linear
  Regression Analysis in Data Science Applications.
\newblock In Proceedings of the 2021 IEEE International Conference on
  Intelligent Systems, Smart and Green Technologies (ICISSGT), Visakhapatnam, India, 13--14 November 2021; pp.
  167--170.
\newblock {https://doi.org/10.1109/ICISSGT52025.2021.00043}. %MDPI: newly added information, please confirm. <-- Answer: yes, we confirm.

\bibitem[Ahmed \em{et~al.}(2021)Ahmed, Mahmud, and Tuya]{Ahmed}
Ahmed, R.; Mahmud, K.H.; Tuya, J.H.
\newblock {A GIS-Based Mathematical Approach for Generating 3D Terrain Model
  from High-Resolution UAV Imageries}.
\newblock {\em J. Geovisualization Spat. Anal.} {\bf 2021},
  {\em 5},~24.
\newblock {https://doi.org/10.1007/s41651-021-00094-7}.

\bibitem[Zhou \em{et~al.}(2021)Zhou, Li, Chi, Yin, Oravecz, Bodovski, Friedman,
  Vrieze, and Chow]{Zhou2021GPS2spaceAO}
Zhou, S.; Li, Y.; Chi, G.; Yin, J.; Oravecz, Z.; Bodovski, Y.; Friedman, N.P.;
  Vrieze, S.I.; Chow, S.M.
\newblock GPS2space: An Open-source Python Library for Spatial Measure
  Extraction from GPS Data.
\newblock {\em J. Behav. Data Sci.} {\bf 2021}, {\em
  1},~127--155.
\newblock {https://doi.org/10.35566/jbds/v1n2/p5}.

\bibitem[Jovanov and Naumoski(2020)]{9254627}
Jovanov, S.; Naumoski, A.
\newblock {A GIS-based Mapping of Mountain Peaks, Waterfalls and Mountain
  Lodges in North Macedonia}.
\newblock In Proceedings of the {2020 4th International Symposium on
  Multidisciplinary Studies and Innovative Technologies (ISMSIT)}, Ankara, Turkey, 22--24 October 2020; pp.
  1--5.
\newblock {https://doi.org/10.1109/ISMSIT50672.2020.9254627}. %MDPI: newly added information, please confirm. <-- Answer: yes, we confirm.

\bibitem[Mogaji \em{et~al.}(2022)Mogaji, Atenidegbe, Adeyemo, and
  Akinmulewo]{Mogaji}
Mogaji, K.A.; Atenidegbe, O.F.; Adeyemo, I.A.; Akinmulewo, K.P.
\newblock {Application of GIS-based PROMETHEE data mining technique to
  geoelectrical-derived parameters for aquifer potentiality assessment in a
  typical hardrock terrain Southwestern Nigeria}.
\newblock {\em Sustain. Water Resour. Manag.} {\bf 2022}, {\em 8},~51.
\newblock {https://doi.org/10.1007/s40899-022-00616-1}.

\bibitem[Ma \em{et~al.}(2020)Ma, Longley, Salmond, and Gao]{Maetal2020}
Ma, X.; Longley, I.; Salmond, J.; Gao, J.
\newblock {PyLUR: Efficient software for land use regression modeling the
  spatial distribution of air pollutants using GDAL/OGR library in Python}.
\newblock {\em Front. Environ. Sci. Eng.} {\bf 2020},
  {\em 14},~44.
\newblock {https://doi.org/10.1007/s11783-020-1221-5}.

\bibitem[Dobesova(2011)]{6057428}
Dobesova, Z.
\newblock {Programming language Python for data processing}.
\newblock In Proceedings of the 2011 International Conference on Electrical and
  Control Engineering, Yichang, China, 16--18 September 2011; pp. 4866--4869.
\newblock {https://doi.org/10.1109/ICECENG.2011.6057428}. %MDPI: newly added information, please confirm. <-- Answer: yes, we confirm.

\bibitem[Jaskolka \em{et~al.}(2019)Jaskolka, Seiler, Beyer, and
  Kaup]{JASKOLKA2019e02560}
Jaskolka, K.; Seiler, J.; Beyer, F.; Kaup, A.
\newblock {A Python-based laboratory course for image and video signal
  processing on embedded systems}.
\newblock {\em Heliyon} {\bf 2019}, {\em 5},~e02560.
\newblock {https://doi.org/10.1016/j.heliyon.2019.e02560}.

\bibitem[Bastidas \em{et~al.}(2021)Bastidas, Auras, and
  Juurlink]{BASTIDAS2021150821}
Bastidas, J.M.P.; Auras, S.V.; Juurlink, L.B.F.
\newblock {A Python script to automate STM image analysis for stepped
  surfaces}.
\newblock {\em Appl. Surf. Sci.} {\bf 2021}, {\em 567},~150821.
\newblock {https://doi.org/10.1016/j.apsusc.2021.150821}.

\bibitem[{R Core Team}(2022)]{RCoreTeam}
{R Core Team}.
\newblock {\em R: A Language and Environment for Statistical Computing};
\newblock R Foundation for Statistical Computing: Vienna, Austria,  2022.

\bibitem[Murrell(2005)]{Murrell}
Murrell, P.
\newblock {\em R Graphics}, 1st ed.; Chapman and Hall/CRC: New York, NY, USA,  2005.
\newblock {https://doi.org/10.1201/9781420035025}.

\bibitem[Chung \em{et~al.}(2016)Chung, Ibrahim, Hassan, and Rosli]{7847824}
Chung, T.D.; Ibrahim, R.; Hassan, S.M.; Rosli, N.S.
\newblock Fast approach for automatic data retrieval using R programming
  language.
\newblock In Proceedings of the 2016 2nd IEEE International Symposium on
  Robotics and Manufacturing Automation (ROMA), Ipoh, Malaysia, 25--27 September 2016; pp. 1--4.
\newblock {https://doi.org/10.1109/ROMA.2016.7847824}.

\bibitem[Lemenkova(2022)]{Lemenkova2022a}
Lemenkova, P.
\newblock {Tanzania Craton, Serengeti Plain and Eastern Rift Valley: mapping of
  geospatial data by scripting techniques}.
\newblock {\em Est. J. Earth Sci.} {\bf 2022}, {\em
  71},~61--79.
\newblock {https://doi.org/10.3176/earth.2022.05}.

\bibitem[Al-Amin \em{et~al.}(2020)Al-Amin, Uday Sampreeth~Chebolu, and
  Ordonez]{9378399}
Al-Amin, S.T.; Uday Sampreeth~Chebolu, S.; Ordonez, C.
\newblock Extending the R Language with a Scalable Matrix Summarization
  Operator.
\newblock In Proceedings of the 2020 IEEE International Conference on Big Data
  (Big Data), Atlanta, GA, USA, 10--13 December 2020; pp. 399--405.
\newblock {https://doi.org/10.1109/BigData50022.2020.9378399}. %MDPI: newly added information, please confirm. <-- Answer: yes, we confirm.

\bibitem[Lemenkova(2022)]{Lemenkova202216}
Lemenkova, P.
\newblock {A Script-Driven Approach to Mapping Satellite-Derived Topography and
  Gravity Data Over the Zagros Fold-and-Thrust Belt, Iran}.
\newblock {\em Artif. Satell.} {\bf 2022}, {\em 57},~110--137.
\newblock {https://doi.org/10.2478/arsa-2022-0006}.

\bibitem[Wang \em{et~al.}(2021)Wang, Wei, and Bai]{9390011}
Wang, D.; Wei, H.; Bai, B.
\newblock Teaching Design and Implementation Based on R Language Under the
  Background of Big Data.
\newblock In Proceedings of the 2021 IEEE 2nd International Conference on Big
  Data, Artificial Intelligence and Internet of Things Engineering (ICBAIE),
  Nanchang, China, 26--28 March  2021; pp. 49--52.
\newblock {https://doi.org/10.1109/ICBAIE52039.2021.9390011}.

\bibitem[Malviya \em{et~al.}(2016)Malviya, Udhani, and Soni]{7570960}
Malviya, A.; Udhani, A.; Soni, S.
\newblock R-tool: Data analytic framework for big data.
\newblock In Proceedings of the 2016 Symposium on Colossal Data Analysis and
  Networking (CDAN),  Indore, India, 18--19 March 2016; pp. 1--5.
\newblock {https://doi.org/10.1109/CDAN.2016.7570960}.

\bibitem[Wang(2022)]{9752002}
Wang, S.
\newblock The Design of Medical English Autonomous Guiding Platform under the
  Information Technology Environment-Based on R Language.
\newblock In Proceedings of the 2022 International Conference on Electronics
  and Renewable Systems (ICEARS), Tuticorin, India, 16-18 March 2022; pp. 1584--1587.
\newblock {https://doi.org/10.1109/ICEARS53579.2022.9752002}. %MDPI: newly added information, please confirm. <-- Answer: yes, we confirm.

\bibitem[Lin \em{et~al.}(2013)Lin, Yang, and Midkiff]{6597171}
Lin, H.; Yang, S.; Midkiff, S.P.
\newblock RABID -- A General Distributed R Processing Framework Targeting Large
  Data-Set Problems.
\newblock In Proceedings of the 2013 IEEE International Congress on Big Data,
 Santa Clara, CA, USA, 27 June--2 July 2013; pp. 423--424.
\newblock {https://doi.org/10.1109/BigData.Congress.2013.67}. %MDPI: newly added information, please confirm. <-- Answer: yes, we confirm.

\bibitem[Wang \em{et~al.}(2017)Wang, Xu, Duan, Zhang, and Xu]{7978404}
Wang, G.; Xu, Y.; Duan, Q.; Zhang, M.; Xu, B.
\newblock Prediction model of glutamic acid production of data mining based on
  R language.
\newblock In Proceedings of the 2017 29th Chinese Control And Decision
  Conference (CCDC), Chongqing, China, 28--30 May 2017; pp.~6806--6810.
\newblock {https://doi.org/10.1109/CCDC.2017.7978404}. %MDPI: newly added information, please confirm. <-- Answer: yes, we confirm.

\bibitem[Lemenkova(2019)]{Lemenkova201991}
Lemenkova, P.
\newblock {Statistical Analysis of the Mariana Trench Geomorphology Using R
  Programming Language}.
\newblock {\em Geod. Cartogr.} {\bf 2019}, {\em 45},~57--84.
\newblock {https://doi.org/10.3846/gac.2019.3785}.

\bibitem[Mao(2021)]{9675988}
Mao, A.
\newblock Construction of Intelligent Vocational Management Information System
  with R Programming.
\newblock In Proceedings of the 2021 5th International Conference on
  Electronics, Communication and Aerospace Technology (ICECA), Coimbatore, India, 1--3 December 2021; pp.
  1162--1165.
\newblock {https://doi.org/10.1109/ICECA52323.2021.9675988}. %MDPI: newly added information, please confirm. <-- Answer: yes, we confirm.

\bibitem[Bishwal(2017)]{8226275}
Bishwal, R.M.
\newblock Potential use of R-statistical programming in the field of
  geoscience.
\newblock In Proceedings of the 2017 2nd International Conference for
  Convergence in Technology (I2CT), Pune, India, 8--9 April 2017; pp. 979--982.
\newblock {https://doi.org/10.1109/I2CT.2017.8226275}.

\bibitem[Frery \em{et~al.}(2022)Frery, Wu, and Gomez]{9844952}
Frery, A.C.; Wu, J.; Gomez, L., Elements of Data Analysis and Image Processing
  with R.
\newblock In {\em SAR Image Analysis---A Computational Statistics Approach:
  With R Code, Data, and Applications}; Willey:  Hoboken, NJ, USA, %MDPI: newly added information, please confirm. <-- Answer: yes, we confirm.
 2022; pp. 43--60.
\newblock {https://doi.org/10.1002/9781119795520.ch2}.

\bibitem[Ramalakshmi and Kompala(2017)]{8117873}
Ramalakshmi, E.; Kompala, N.
\newblock Multi-threading image processing in single-core and multi-core CPU
  using R language.
\newblock In Proceedings of the 2017 Second International Conference on
  Electrical, Computer and Communication Technologies (ICECCT), Erode, India, 22--24 February 2017; pp.
  1--5.
\newblock {https://doi.org/10.1109/ICECCT.2017.8117873}.

\bibitem[Chatelain \em{et~al.}(2010)Chatelain, Bakayoko, Martin, and
  Gautier]{Chatelainetal2010}
Chatelain, C.; Bakayoko, A.; Martin, P.; Gautier, L.
\newblock {Monitoring tropical forest fragmentation in the Zagné-Taï area
  (west of Taï National Park, Côte d’Ivoire)}.
\newblock {\em Biodivers. Conserv.} {\bf 2010}, {\em 19},~2405--2420.
\newblock {https://doi.org/10.1007/s10531-010-9847-4}.

\bibitem[Hennenberg \em{et~al.}(2008)Hennenberg, Orthmann, Steinke, and
  Porembski]{Hennenberg}
Hennenberg, K.J.; Orthmann, B.; Steinke, I.; Porembski, S.
\newblock {Core area analysis at semi-deciduous forest islands in the Comoé
  National Park, NE Ivory Coast}.
\newblock {\em Biodivers. Conserv. Vol.} {\bf 2008}, {\em
  17},~2787--2797.
\newblock {https://doi.org/10.1007/s10531-007-9292-1}.

\bibitem[El-Shahat \em{et~al.}(2021)El-Shahat, El-Zafarany, El~Seoud, and
  Ghoniem]{ElShahat}
El-Shahat, S.; El-Zafarany, A.M.; El~Seoud, T.A.; Ghoniem, S.A.
\newblock {Vulnerability assessment of African coasts to sea level rise using
  GIS and remote sensing}.
\newblock {\em Environ. Dev. Sustain.} {\bf 2021}, {\em
  23},~2827--2845.
\newblock {https://doi.org/10.1007/s10668-020-00639-8}.

\bibitem[Tang \em{et~al.}(2016)Tang, Feng, Jia, Shi, Zuo, and
  Trettin]{Tangetal2016}
Tang, W.; Feng, W.; Jia, M.; Shi, J.; Zuo, H.; Trettin, C.C.
\newblock {The assessment of mangrove biomass and carbon in West Africa: a
  spatially explicit analytical framework}.
\newblock {\em Wetl. Ecol. Manag.} {\bf 2016}, {\em 24},~153--171.
\newblock {https://doi.org/10.1007/s11273-015-9474-7}.

\bibitem[Affian \em{et~al.}(2009)Affian, Robin, Maanan, Digbehi, Djagoua, and
  Kouamé]{Affianetal2009}
Affian, K.; Robin, M.; Maanan, M.; Digbehi, B.; Djagoua, E.V.; Kouamé, F.
\newblock {Heavy metal and polycyclic aromatic hydrocarbons in Ebrié lagoon
  sediments, Côte d’Ivoire}.
\newblock {\em Environ. Monit. Assess.} {\bf 2009}, {\em
  159},~531.
\newblock {https://doi.org/10.1007/s10661-008-0649-z}.

\bibitem[Msiska(2009)]{MSISKA2009487}
Msiska, O.V.
\newblock {Lakes and Reservoirs of Africa: South of Sahara}. In {\em
  Encyclopedia of Inland Waters}; Likens, G.E., Ed.; Academic Press: Oxford, UK,
  2009; pp. 487--500.
\newblock {https://doi.org/10.1016/B978-012370626-3.00036-3}.

\bibitem[Thomasset(1900)]{Thomasset}
Thomasset.
\newblock {La Côte d'Ivoire}.
\newblock {\em Ann. GéOgraphie} {\bf 1900}, {\em 9},~159--172.
\newblock {https://doi.org/10.3406/geo.1900.6213}.

\bibitem[Rougerie(1959)]{Rougerie}
Rougerie, G.
\newblock {Façonnement actuel des modèles en Cote d'Ivoire forestière}.
\newblock {\em L'Inform. Géograph.} {\bf 1959}, {\em 23},~135--136.

\bibitem[Tricart(1957)]{Tricart}
Tricart, J.
\newblock {Le café en Côte d'Ivoire}.
\newblock {\em Cah. D'Outre-Mer} {\bf 1957}, {\em 39},~209--233.
\newblock {https://doi.org/10.3406/caoum.1957.2042}.

\bibitem[Bruno(2000)]{Losch}
Bruno, L.
\newblock {Coup de cacao en Côte d'Ivoire}.
\newblock {\em Crit. Int.} {\bf 2000}, {\em 9},~6--14.
\newblock {https://doi.org/10.3406/criti.2000.1614}.

\bibitem[{de Planhol}(1947)]{Planhol}
{De Planhol}, X.
\newblock {Le cacao en Côte d’Ivoire: Étude de géographie régionale}.
\newblock {\em L'Inform. Géograph.} {\bf 1947}, {\em 11},~50--57.

\bibitem[{Smith Dumont} \em{et~al.}(2014){Smith Dumont}, Gnahoua, Ohouo,
  Sinclair, and Vaast]{Smithetal2014}
{Smith Dumont}, E.; Gnahoua, G.M.; Ohouo, L.; Sinclair, F.L.; Vaast, P.
\newblock {Farmers in Côte d’Ivoire value integrating tree diversity in
  cocoa for the provision of ecosystem services}.
\newblock {\em Agrofor. Syst.} {\bf 2014}, {\em 88},~1047--1066.
\newblock {https://doi.org/10.1007/s10457-014-9679-4}.

\bibitem[Läderach \em{et~al.}(2013)Läderach, Martinez-Valle, Schroth, and
  Castro]{Laderach}
Läderach, P.; Martinez-Valle, A.; Schroth, G.; Castro, N.
\newblock {Predicting the future climatic suitability for cocoa farming of the
  world’s leading producer countries, Ghana and Côte d’Ivoire}.
\newblock {\em Clim. Chang.} {\bf 2013}, {\em 119},~841--854.
\newblock {https://doi.org/10.1007/s10584-013-0774-8}.

\bibitem[Sawadogo(1974)]{Sawadogo}
Sawadogo, A.
\newblock {La stratégie du développement de l'agriculture en Côte-d'Ivoire}.
\newblock {\em Bull. L'Assoc. Géograph. Français} {\bf 1974},
  {\em 415-416},~87--103.
\newblock {https://doi.org/10.3406/bagf.1974.4760}.

\bibitem[Pélissier(1974)]{Pelissier}
Pélissier, P.
\newblock {Agriculture et développement l'exemple de la Côte-d'Ivoire}.
\newblock {\em Bull. L'Assoc. Géograph. Français} {\bf 1974},
  {\em 415-416},~81--85.
\newblock {https://doi.org/10.3406/bagf.1974.4759}.

\bibitem[Sako \em{et~al.}(2018)Sako, Semdé, and Wenmenga]{SAKO2018297}
Sako, A.; Semdé, S.; Wenmenga, U.
\newblock Geochemical evaluation of soil, surface water and groundwater around
  the Tongon gold mining area, northern Côte d’Ivoire, West Africa.
\newblock {\em J. Afr. Earth Sci.} {\bf 2018}, {\em
  145},~297--316.
\newblock {https://doi.org/10.1016/j.jafrearsci.2018.05.016}.

\bibitem[Sauerwein(2020)]{SAUERWEIN2020104323}
Sauerwein, T.
\newblock Gold mining and development in Côte d’Ivoire: Trajectories,
  opportunities and oversights.
\newblock {\em Land Use Policy} {\bf 2020}, {\em 91},~104323.
\newblock {https://doi.org/10.1016/j.landusepol.2019.104323}.

\bibitem[Murray \em{et~al.}(2019)Murray, Torvela, and Bills]{MURRAY2019351}
Murray, S.; Torvela, T.; Bills, H.
\newblock A geostatistical approach to analyzing gold distribution controlled
  by large-scale fault systems---An example from Côte d’Ivoire.
\newblock {\em J. Afr. Earth Sci.} {\bf 2019}, {\em
  151},~351--370.
\newblock {https://doi.org/10.1016/j.jafrearsci.2018.12.019}.

\bibitem[Sournia(2003)]{Sournia}
Sournia, G.
\newblock {Aménagement du territoire et stratégie du développement en
  Côte-d'Ivoire}.
\newblock {\em L'information Géographique} {\bf 2003}, {\em 67},~124--129.
\newblock {https://doi.org/10.3406/ingeo.2003.2862}.

\bibitem[Chatelain \em{et~al.}(1996)Chatelain, Gautier, and
  Spichiger]{Chatelain}
Chatelain, C.; Gautier, L.; Spichiger, R.
\newblock {A recent history of forest fragmentation in southwestern Ivory
  Coast}.
\newblock {\em Biodivers. Conserv.} {\bf 1996}, {\em 5},~37--53.
\newblock {https://doi.org/10.1007/BF00056291}.

\bibitem[{Denguéadhé Kolongo} \em{et~al.}(2006){Denguéadhé Kolongo},
  Decocq, Yao, Blom, and {Van Rompaey}]{Kolongo}
{Denguéadhé Kolongo}, T.S.; Decocq, G.; Yao, C.Y.A.; Blom, E.C.; {Van
  Rompaey}, R.S.A.R.
\newblock {Plant Species Diversity in the Southern Part of the Taï National
  Park (Côte d’Ivoire)}.
\newblock {\em Biodivers. Conserv.} {\bf 2006}, {\em 15},~2123--2142.
\newblock {https://doi.org/10.1007/s10531-004-6686-1}.

\bibitem[Yeo \em{et~al.}(2017)Yeo, Delsinne, Konate, Alonso, Aïdara, and
  Peeters]{Yeo}
Yeo, K.; Delsinne, T.; Konate, S.; Alonso, L.L.; Aïdara, D.; Peeters, C.
\newblock {Diversity and distribution of ant assemblages above and below ground
  in a West African forest-savannah mosaic (Lamto, Côte d’Ivoire)}.
\newblock {\em Insectes Sociaux} {\bf 2017}, {\em 64},~155--168.
\newblock {https://doi.org/10.1007/s00040-016-0527-6}.

\bibitem[Dubresson(2003)]{Dubresson}
Dubresson, A.
\newblock {Industrialisation et urbanisation en Côte-d'Ivoire. Contribution
  géographique à l'étude de l'accumulation urbaine}.
\newblock {\em L'Inform. Géograph.} {\bf 2003}, {\em 67},~130--133.
\newblock {https://doi.org/10.3406/ingeo.2003.2863}.

\bibitem[Chaléard(1995)]{Chaleard}
Chaléard, J.L.
\newblock {Temps des villes, temps des vivres. L'essor du vivrier marchand en
  Côte d'Ivoire}.
\newblock {\em L'Inform. Géograph.} {\bf 1995}, {\em 59},~42--43.

\bibitem[Cotten(1974)]{Cotten1974}
\textls[-35]{Cotten, A.M.
\newblock {Un aspect de l'urbanisation en Côte-d'Ivoire}.
\newblock {\em Cahiers D'outre-mer} {\bf 1974}, {\em 106},~183--193.
\newblock {https://doi.org/10.3406/caoum.1974.2697}.}

\bibitem[Cotten(1973)]{Cotten1973}
Cotten, A.M.
\newblock {Le rôle des villes moyennes en Côte-d'Ivoire}.
\newblock {\em Bull. L'Inform. Géograph. Français} {\bf 1973},
  {\em 410},~619--625.
\newblock {https://doi.org/10.3406/bagf.1973.4719}.

\bibitem[Cotten(1985)]{Cotten1985}
Cotten, A.M.
\newblock {Développement des transports en République de Côte d'Ivoire Ses
  conséquences géographiques}.
\newblock {\em Trav.  l'Inst. Géograph. Reims} {\bf 1985}, {\em
  63-64},~85--94.
\newblock 
  {https://doi.org/10.3406/tigr.1985.1176}.

\bibitem[Doumouya \em{et~al.}(2012)Doumouya, Dibi, Kouame, Saley, Jourda,
  Savane, and Biemi]{Doumouya}
Doumouya, I.; Dibi, B.; Kouame, K.I.; Saley, B.; Jourda, J.P.; Savane, I.;
  Biemi, J.
\newblock {Modelling of favourable zones for the establishment of water points
  by geographical information system (GIS) and multicriteria analysis (MCA) in
  the Aboisso area (South-east of Côte d’Ivoire)}.
\newblock {\em Environ. Earth Sci.} {\bf 2012}, {\em 67},~1763--1780.
\newblock {https://doi.org/10.1007/s12665-012-1622-2}.

\bibitem[{U.S. Geological Survey}(2015)]{USGS}
{U.S. Geological Survey}.
\newblock {\em Landsat---Earth Observation Satellites};
\newblock Technical Report; USGS: Asheville, NC, USA,  2015.
\newblock {https://doi.org/10.3133/fs20153081}.

\bibitem[{Department of the Interior U.S. Geological Survey}(2022)]{Landsat9}
{Department of the Interior U.S. Geological Survey}.
\newblock {\em Landsat 9 Data Users Handbook};
\newblock EROS: Sioux Falls, SD, USA,  2022;
\newblock LSDS-2082 Version 1.0.

\bibitem[{GEBCO Compilation Group}(2020)]{GEBCO}
{GEBCO Compilation Group}.
\newblock {GEBCO 2020 Grid}. 2020.
\newblock Available online: \url{https://doi.org/10.5285/a29c5465-b138-234d-e053-6c86abc040b9} (accessed on 23 November 2022).%MDPI: Please add the access date (format: Date Month Year), e.g., accessed on 1 January 2020. <-- Answer: added.

\bibitem[Wessel \em{et~al.}(2019)Wessel, Luis, Uieda, Scharroo, Wobbe, Smith,
  and Tian]{Wessel}
Wessel, P.; Luis, J.F.; Uieda, L.; Scharroo, R.; Wobbe, F.; Smith, W.H.F.;
  Tian, D.
\newblock {The Generic Mapping Tools version 6}.
\newblock {\em Geochem. Geophys. Geosyst.} {\bf 2019}, {\em
  20},~5556--5564.
\newblock {xhttps://doi.org/10.1029/2019GC008515}.

\bibitem[Lemenkova(2022{\natexlab{a}})]{Lemenkova2022b}
Lemenkova, P.
\newblock {Console-Based Mapping of Mongolia Using GMT Cartographic Scripting
  Toolset for Processing TerraClimate Data}.
\newblock {\em Geosciences} {\bf 2022}, {\em 12},~140.
\newblock {https://doi.org/10.3390/geosciences12030140}.

\bibitem[Lemenkova(2022{\natexlab{b}})]{Lemenkova202217}
Lemenkova, P.
\newblock {Mapping Climate Parameters over the Territory of Botswana Using GMT
  and Gridded Surface Data from TerraClimate}.
\newblock {\em ISPRS Int. J. Geo. Inf.} {\bf 2022}, {\em
  11},~473.
\newblock {https://doi.org/10.3390/ijgi11090473}.

\bibitem[{R. J. Hijmans and R. Bivand and K. Forner and J. Ooms and E. Pebesma
  and M. D. Sumner}(2022)]{RTerra}
{ Hijmans, R.J.; Bivand, R.; Forner, K.;  Ooms, J.; Pebesma, E.; 
  Sumner, M.D.}
\newblock {\em Package ‘Terra’};
\newblock Maintainer: Vienna, Austria,  2022.

\bibitem[Hijmans(2017)]{Hijmans}
Hijmans, R.J.
\newblock Raster: Geographic Data Analysis and Modeling. R Package Version
  2.6-7. 2017.
\newblock Available online:
\newblock \url{https://CRAN.R-project.org/package=raster} (accessed on 23 November 2022).%MDPI: Please add the access date (format: Date Month Year), e.g., accessed on 1 January 2020. <-- Answer: added.

\bibitem[Neuwirth(2014)]{Neuwirth}
Neuwirth, E.
\newblock RColorBrewer: ColorBrewer Palettes. R Package Version 1.1-2. 2014.
\newblock Available online:  
\newblock \url{https://CRAN.R-project.org/package=RColorBrewer} (accessed on 23 November 2022).%MDPI: Please add the access date (format: Date Month Year), e.g., accessed on 1 January 2020. <-- Answer: added.

\bibitem[Volesky \em{et~al.}(2003)Volesky, Stern, and Johnson]{VOLESKY2003235}
Volesky, J.C.; Stern, R.J.; Johnson, P.R.
\newblock Geological control of massive sulfide mineralization in the
  Neoproterozoic Wadi Bidah shear zone, southwestern Saudi Arabia, inferences
  from orbital remote sensing and field studies.
\newblock {\em Precambrian Res.} {\bf 2003}, {\em 123},~235--247.
\newblock {https://doi.org/10.1016/S0301-9268(03)00070-6}. %MDPI: Please check if this could be removed. <-- Answer: yes, agree (removed).

\bibitem[Richards(2013)]{Richards}
Richards, J.A.
\newblock {\em Remote Sensing Digital Image Analysis. An Introduction}, 5th ed.;
  Springer: Dordrecht, The Netherlands, 2013.
\newblock {https://doi.org/10.1007/978-3-642-30062-2}.

\bibitem[Campbell and Wynne(2011)]{CampbellWynne}
Campbell, J.B.; Wynne, R.H.
\newblock {\em Introduction to Remote Sensing}, 5th ed.; Guilford Press: New
  York, NY, USA, 2011.

\bibitem[Deoli \em{et~al.}(2022)Deoli, Kumar, and Kuriqi]{s22186827}
Deoli, V.; Kumar, D.; Kuriqi, A.
\newblock Detection of Water Spread Area Changes in Eutrophic Lake Using
  Landsat Data.
\newblock {\em Sensors} {\bf 2022}, {\em 22}.
\newblock {https://doi.org/10.3390/s22186827}.

\bibitem[Hunter(2007)]{Hunter2007MatplotlibA2}
Hunter, J.D.
\newblock Matplotlib: A 2D Graphics Environment.
\newblock {\em Comput. Sci. Eng.} {\bf 2007}, {\em
  9},~90--95.
\newblock {https://doi.org/10.1109/MCSE.2007.55}.

\bibitem[Gillies \em{et~al.}(2013)Gillies et~al.]{gillies_2019}
Gillies, S. 
\newblock Rasterio: geospatial raster I/O for {Python} programmers, 2013
\newblock Available online:  
\newblock \url{https://rasterio.readthedocs.io/en/latest/} (accessed on 23 November 2022)
%MDPI: Please provide more information about the article type, such as book (please provide the name and location of the publisher); online resource (please provide the URL of the website and the date it was accessed (Date Month Year)); or journal article (please provide the name of the journal, the year and volume in which it was published, and the page number). Please refer to https://www.mdpi.com/authors/references for full reference formatting guides. <--  Answer: added information about the online resource with URL and access date.
.

\bibitem[Harris \em{et~al.}(2020)Harris, Millman, van~der Walt, Gommers,
  Virtanen, Cournapeau, Wieser, Taylor, Berg, Smith, Kern, Picus, Hoyer, van
  Kerkwijk, Brett, Haldane, del R{\'{i}}o, Wiebe, Peterson,
  G{\'{e}}rard-Marchant, Sheppard, Reddy, Weckesser, Abbasi, Gohlke, and
  Oliphant]{harris2020array}
Harris, C.R.; Millman, K.J.; van~der Walt, S.J.; Gommers, R.; Virtanen, P.;
  Cournapeau, D.; Wieser, E.; Taylor, J.; Berg, S.; Smith, N.J.;  et~al.
\newblock Array programming with {NumPy}.
\newblock {\em Nature} {\bf 2020}, {\em 585},~357--362.
\newblock {https://doi.org/10.1038/s41586-020-2649-2}.

\bibitem[{Plotly Technologies Inc}(2015)]{plotly}
Plotly Technologies Inc.
\newblock {\em Collaborative Data Science}; Plotly Technologies Inc.: Montreal, QC, USA, 2015.

\bibitem[{Rouse} \em{et~al.}(1974){Rouse}, {Haas}, {Schell}, and
  {Deering}]{1974NASSP}
{Rouse}, J.~W., J.; {Haas}, R.H.; {Schell}, J.A.; {Deering}, D.W.
\newblock {Monitoring Vegetation Systems in the Great Plains with Erts}. In
  {\em NASA Special Publication}; NASA: Washington, DC, USA, 1974; Volume 351, p. 309.

\bibitem[Huete(1988)]{HUETE1988295}
Huete, A.R.
\newblock A soil-adjusted vegetation index (SAVI).
\newblock {\em Remote Sens. Environ.} {\bf 1988}, {\em 25},~295--309.
\newblock {https://doi.org/10.1016/0034-4257(88)90106-X}.

\bibitem[Jiang \em{et~al.}(2008)Jiang, Huete, Didan, and Miura]{JIANG20083833}
Jiang, Z.; Huete, A.R.; Didan, K.; Miura, T.
\newblock Development of a two-band enhanced vegetation index without a blue
  band.
\newblock {\em Remote Sens. Environ.} {\bf 2008}, {\em
  112},~3833--3845.
\newblock {https://doi.org/10.1016/j.rse.2008.06.006}.

\bibitem[Kaufman and Tanre(1992)]{134076}
Kaufman, Y.J.; Tanre, D.
\newblock Atmospherically resistant vegetation index (ARVI) for EOS-MODIS.
\newblock {\em IEEE Trans. Geosci. Remote. Sens.} {\bf 1992},
  {\em 30},~261--270.
\newblock {https://doi.org/10.1109/36.134076}.

\bibitem[Avand \em{et~al.}(2022)Avand, Kuriqi, Khazaei, and
  Ghorbanzadeh]{AVAND20221}
Avand, M.; Kuriqi, A.; Khazaei, M.; Ghorbanzadeh, O.
\newblock DEM resolution effects on machine learning performance for flood
  probability mapping.
\newblock {\em J. Hydro. Environ. Res.} {\bf 2022}, {\em
  40},~1--16.
\newblock {https://doi.org/10.1016/j.jher.2021.10.002}.

\end{thebibliography}

%\nocite{*}

%\section*{\hl{Short Biography of Authors:} %MDPI: this journal has no this part, please confirm}
%\bio
%{\raisebox{-0.35cm}{\includegraphics[width=3.5cm,height=5.3cm,clip,keepaspectratio]{Photo_author_1.jpg}}}
%{\textbf{Polina Lemenkova} is an Engineer and a PhD student of Prof. Dr. O. Debeir in the Image Analysis Research Unit (LISA-IA) in École polytechnique de Bruxelles, Université Libre de Bruxelles. Her research interests are focused on image processing, cartography, spatial data analysis, remote sensing and applied programming. She has a strong research experience in Earth and environmental sciences, mapping and data visualization, specifically in geological and geophysical domains. Her current topic is on seismic data processing using machine learning methods and scripting languages.}
%%
%\bio
%{\raisebox{-0.35cm}{\includegraphics[width=3.5cm,height=5.3cm,clip,keepaspectratio]{Photo_author_2.jpg}}}
%{\textbf{Olivier Debeir} is a Head of the Image Analysis Research Unit (LISA-IA) in the École polytechnique de Bruxelles (Brussels Faculty of Engineering), Université Libre de Bruxelles. His research interest is in addressing questions related to artificial intelligence and machine learning in general, and specifically on computer vision, graphics, image analysis, segmentation and classification, pattern recognition and deep learning. His research activity also concerns applied statistics and diverse algorithms of data analysis and modelling in natural sciences. He is an expert in programming languages.}

%%%%%%%%%%%%%%%%%%%%%%%%%%%%%%%%%%%%%%%%%%
\end{adjustwidth}
\end{document}

